\documentclass[11pt]{article}

\usepackage[margin=1in]{geometry}
\usepackage{amsmath,amssymb,amsthm,mathtools}
\usepackage{enumitem}
\usepackage{hyperref}
\hypersetup{colorlinks=true,linkcolor=blue,citecolor=blue,urlcolor=blue}
\usepackage{orcidlink}
\usepackage{fontawesome5}

\newtheorem{theorem}{Theorem}[section]
\newtheorem{proposition}[theorem]{Proposition}
\newtheorem{lemma}[theorem]{Lemma}
\newtheorem{corollary}[theorem]{Corollary}
\newtheorem{conjecture}[theorem]{Conjecture}
\newtheorem{empirical}[theorem]{Empirical Result}
\theoremstyle{definition}
\newtheorem{definition}[theorem]{Definition}
\theoremstyle{remark}
\newtheorem{remark}[theorem]{Remark}
\newtheorem{example}[theorem]{Example}

\DeclareMathOperator{\ord}{ord}
\DeclareMathOperator{\Cov}{Cov}
\DeclareMathOperator{\Var}{Var}
\DeclareMathOperator{\Corr}{Corr}
\newcommand{\ZZ}{\mathbb{Z}}
\newcommand{\NN}{\mathbb{N}}
\newcommand{\RR}{\mathbb{R}}
\newcommand{\CC}{\mathbb{C}}
\newcommand{\Zp}{\mathbb{Z}_3}
\newcommand{\Zpu}{\mathbb{Z}_3^{\times}}
\newcommand{\e}[1]{e(#1)}

\title{The $qx+1$ Generalization of Tao's Syracuse Measure:\\
A Closed-Form Pressure Equation, an Endogeny Barrier,\\
and Complementary Viewpoints on the Arithmetic Obstruction}

\author{Renato Augusto Tavares~{\small\href{mailto:dr.renatotavares@gmail.com}{\textcolor{black}{\faEnvelope}}}~\orcidlink{0009-0002-0196-3311}\\
\normalsize Universidade Federal de Goiás\\
\small \href{https://orcid.org/0009-0002-0196-3311}{https://orcid.org/0009-0002-0196-3311}}
\date{\today}

\begin{document}

\maketitle

\begin{abstract}
We study a natural $qx+1$ generalization of Tao's Syracuse random
variable \cite{Tao2022}, the probabilistic object underlying the
strongest unconditional partial result toward the Collatz conjecture.
We prove an exact identity, $\rho_{\mathrm{ann}}(\alpha) =
q^{\alpha-1}/(2^\alpha-1)$, for the \emph{annealed} pressure of the
natural multiplicative population functional on the reverse tree of
the accelerated $qx+1$ map --- the exponential growth rate, averaged
over all root residues at a given depth, of the associated partition
function --- via an exact fibre-counting bijection, correcting a
finite-automaton argument for the same identity that we show is not
well posed. The nontrivial root of $\rho_{\mathrm{ann}}(\alpha)=1$
undergoes a structural transition at $q=5$: it equals exactly
$\alpha^\ast=2$ at $q=3$ and falls strictly below $1$ for every
$q\ge5$, conjecturally signalling a transition from exponent one to a
strictly sublinear counting exponent. For $q=3$ this is the classical
Growth Exponent Conjecture of Kontorovich--Lagarias and
Applegate--Lagarias. We show,
moreover, that the two roots of this equation always sit on opposite
sides of a quenched/annealed freezing transition in the sense of
directed polymers in a random environment: the smaller root is always
unfrozen, so the annealed identity transfers rigorously to the
quenched counting function of the matching i.i.d.\
branching-random-walk model, and, we conjecture on strong empirical
grounds, to the genuine arithmetic tree itself; the larger root is
always frozen, so the pressure identity alone does not establish the
tail index of the density martingale. We prove that the corresponding
tilted $q$-adic measure is singular at every frozen root by an entropy
identity and a fractional-moment estimate. We prove the tail index for the
matching i.i.d. model by an implicit renewal theorem. For the coherent
$q$-adic environment, we prove martingale convergence and identify the
limit with the density of the absolutely continuous component; absolute
continuity and the tail remain conjectural. At $q=3$ the predicted index
$\alpha^\ast=2$ matches three empirical measurements in earlier work
on the classical Collatz map; for $q\ge5$ the arithmetic claim is
strengthened by a large-sample
statistical test at $q=5$ (though still not proved), and still only
weakly supported for $q\ge7$. We then identify and
formally characterize an \emph{endogeny barrier}, in the technical
sense of Aldous--Bandyopadhyay \cite{AldousBandyopadhyay2005},
obstructing any proof that relies only on the linear recursive
structure of the tree: solutions of the form $G = W\cdot Y$, with $Y$
an arbitrary bounded factor invariant along the tree, solve the same
fixed-point recursion as the canonical (``endogenous'') solution $W$
and are statistically indistinguishable from it by every observable we
are able to test. We study the gap through four complementary
vocabularies (endogeny itself counting as the first). Closing
this gap leads to a small collection of unresolved arithmetic
statements: Wirsching's 1998 Weak Covering Conjecture
\cite{Wirsching1998}, Tao's weighted $\beta=1$ conjecture
\cite{Tao2020}, a functional equation for an explicit base-$3$
analogue of the Fabius function arising in Wirsching's 2003 paper
\cite{Wirsching2003}, and two candidate lemmas suggested as natural
next steps by Tao's bivariate Fourier-decay technique
\cite{Tao2022,Tao2011}, whose primitive Fourier energy is the exact
level increment of the $L^2$ collision mass. A diagonalization in
discrete-logarithm coordinates proves exponential collision growth for
every odd prime $q\ge5$ whose order of $2$ lifts maximally, and reduces
the critical case $q=3$ to an explicit signed spectral balance. We
express the Syracuse law exactly as a geometric
mixture of Wirsching's fixed-cost generators. This proves that his
target uniform multiplicity condition implies Tao's $\beta=1$
cylinder estimate. The first Fourier route
yields a genuine closed-form structure theorem together with a
conditional reduction whose hypothesis (square-integrability of the
Syracuse density on $\mathbb{Z}_3$) is tested by an exact finite-level
recursion and shows sustained linear growth through level $17$. The
other leaves a frequency-stability step open and gives
an exact Jensen calculation showing slack by a factor $\approx1.88$ in
one annealed $\ell^1$ Fourier budget. We complement this
theoretical package with a fully independent, purely empirical result:
implementing the exact reverse tree of the $5x+1$ map and applying
Aitken $\Delta^2$ extrapolation to a sequence of nested-window measurements,
we compare two competing predictions for
the $5x+1$ counting exponent (Kontorovich--Lagarias
\cite{KontorovichLagarias2009} versus a competing stochastic model due
to Volkov, both discussed in \cite{KontorovichLagarias2009}), measuring
$0.639$ (bootstrap interval $[0.633,0.645]$). The remaining monotone
finite-window drift prevents this interval from deciding between the
two asymptotic predictions. Finally, a directed literature search situates
the barrier precisely within the wider theory of Fourier decay for
self-similar and self-conformal measures \cite{BakerBanaji2026,
Bremont}, showing why three standard mechanisms do not directly apply
to this rigid $3$-adic system,
and compares the resulting obstruction with a recent fixed-modulus
orbit statistic \cite{Chang2026}. We distinguish exact identities,
conditional reductions, and analogies, and state the remaining
estimates separately.
\end{abstract}

\noindent\textbf{Keywords:} Collatz conjecture; $qx+1$ map; Syracuse random variable;
pressure equation; freezing transition; branching random walk; martingale tail index;
endogeny; Weak Covering Conjecture; Fabius function; Fourier decay; counting exponent.

\tableofcontents

\section{Introduction}

\subsection{The Collatz conjecture and the Syracuse random variable}

The Collatz (or $3x+1$) conjecture asserts that iterating the map
\[
  C(n) = \begin{cases} n/2 & n \text{ even} \\ 3n+1 & n \text{ odd} \end{cases}
\]
started at any positive integer $n$ eventually reaches $1$. Despite its
elementary statement, the conjecture remains open, and the strongest
unconditional result toward it is due to Tao \cite{Tao2022}, who showed
that \emph{almost all} orbits (in the sense of logarithmic density)
attain almost bounded values. Tao's proof introduces the
\emph{Syracuse random variable}: fixing $n$, let $a_1,a_2,\dots,a_n$ be
i.i.d.\ geometric random variables with $\Pr(a_i=k)=2^{-k}$ for
$k\ge 1$, and define
\begin{equation}\label{eq:syracuse-def}
  F_n(a) := \sum_{j=1}^n 3^{j-1}\, 2^{-a_{[1,j]}} \pmod{3^n}, \qquad
  a_{[1,j]} := a_1+\cdots+a_j,
\end{equation}
where $2^{-1}$ denotes the multiplicative inverse of $2$ modulo $3^n$.
The random variable $\mathrm{Syrac}(\mathbb{Z}/3^n\mathbb{Z}) :=
F_n(a_1,\dots,a_n)$ models, in a precise sense, the residue of an
accelerated Collatz orbit modulo $3^n$ after $n$ odd steps, and Tao's
proof reduces the conjecture's target statement to controlling the
distribution of this random variable, in particular via a Fourier
decay estimate:

\begin{proposition}[Tao 2022, Prop.~1.17]\label{prop:tao117}
For every $A>0$ there is $C_A$ such that for all $n\ge 1$ and every
$\xi\in\mathbb{Z}/3^n\mathbb{Z}$ with $3\nmid \xi$,
\[
  \bigl|\mathbb{E}\, e\bigl(\xi\, F_n(a)/3^n\bigr)\bigr| \le C_A\, n^{-A},
\]
with $C_A$ uniform in $n$ and $\xi$ (subject only to $3\nmid\xi$).
\end{proposition}

This is a super-polynomial (though not exponential) decay estimate,
and it is the technical heart of Tao's theorem. A recursive
consequence of the construction \eqref{eq:syracuse-def}, which we use
repeatedly below, is the following \emph{exact self-similarity}
identity, verified directly against Tao's text:

\begin{proposition}[Tao 2022, eq.~(1.23)]\label{prop:tao123}
For $k\le n$,
\[
  \mathrm{Syrac}(\mathbb{Z}/3^n\mathbb{Z}) \bmod 3^k \; \equiv\;
  \mathrm{Syrac}(\mathbb{Z}/3^k\mathbb{Z})
\]
as an identity in law. Equivalently, writing $F_n(a) =
2^{-a_1}\bigl(1+3\,F_{n-1}(a_2,\dots,a_n)\bigr) \bmod 3^n$, the
coarsest $k$ coordinates of $\mathrm{Syrac}(\mathbb{Z}/3^n\mathbb{Z})$
depend only on $(a_1,\dots,a_k)$, the digits \emph{adjacent to the
root} in the indexing of \eqref{eq:syracuse-def}, not on the deepest
ones.
\end{proposition}

The orientation convention embedded in Proposition~\ref{prop:tao123}
(that $a_1$, not $a_n$, is the coordinate adjacent to the root) is a
recurring source of error when translating Tao's construction into
reverse-tree language, and we flag it explicitly throughout.

\subsection{This paper's contributions}

Building on \cite{Tao2022} and on the classical predecessor-density
program of Wirsching \cite{Wirsching1998,Wirsching2003}, we make the
following contributions.

\begin{enumerate}[label=(\roman*)]
\item \textbf{A closed-form pressure equation for $qx+1$
  (\S\ref{sec:pressure}).} We show that the tail exponent controlling
  the reverse-tree population martingale of the accelerated $qx+1$ map
  satisfies the exact equation $q^{\alpha-1}=2^\alpha-1$, with a
  structural transition at $q=5$.
\item \textbf{A precise endogeny barrier (\S\ref{sec:endogeny}).} We
  exhibit an explicit family of \emph{non-endogenous} solutions to the
  tree's fixed-point recursion, prove they are statistically
  indistinguishable from the canonical one by pressure, tail index, or
  linear regression against the limiting measure. We identify the
  missing arithmetic input as cancellation of centered fine-frequency
  correlations after averaging over path pairs; pairwise independence
  is impossible.
\item \textbf{An empirical bound on the residual coupling
  (\S\ref{sec:calibration}).} A three-arm control experiment gives
  $\rho_{\mathrm{eff}} \lesssim 0.06$ (95\% CI) for the effective
  arithmetic coupling between sibling subtrees, converting the
  theoretical gap of (ii) into a quantitative empirical bound.
\item \textbf{A finite-window empirical result (\S\ref{sec:kl}).} Exact
  enumeration and Aitken $\Delta^2$ extrapolation compare two predictions
  for the $5x+1$ counting exponent, while retaining a visible
  pre-asymptotic window bias.
\item \textbf{A direct computational test of the Weak Covering
  Conjecture (\S\ref{sec:wcc}).}
\item \textbf{Three convergent regimes of precision for Tao's bivariate
  extension (\S\ref{sec:regimes}), and triangulation with two further
  complementary formulations of the arithmetic obstruction
  (\S\ref{sec:triangulation}).}
\item \textbf{Two candidate routes (\S\ref{sec:lemmas}).} The first
  yields a structural obstruction and a finite-level test of a natural
  sufficient condition.
  The second leaves an explicit frequency-stability step open and gives
  an exact Jensen calculation for one annealed Fourier budget.
\item \textbf{A literature-theoretic contextualization
  (\S\ref{sec:context}).} We show why three standard methods from the
  theory of Fourier decay for self-similar measures do not transfer
  directly. We also test a recent fixed-modulus orbit statistic
  \cite{Chang2026}, without treating it as a proved reduction.
\end{enumerate}

We deliberately frame this paper not as an attempted proof of any part
of the Collatz conjecture, but as a precise characterization of an
obstruction that recurs, in essentially the same form, across every
technique we or the existing literature have brought to bear on this
specific line of attack. We return to this framing choice in
\S\ref{sec:discussion}.

\section{Setup: the accelerated $qx+1$ map and its reverse tree}
\label{sec:setup}

Fix an odd integer $q\ge 3$ coprime to $2$. The \emph{accelerated
$qx+1$ map} is
\[
  T_q(n) = \frac{qn+1}{2^{v_2(qn+1)}}, \qquad n \text{ odd},
\]
where $v_2$ is the $2$-adic valuation. For $q=3$ this is the standard
accelerated Collatz map. The \emph{reverse tree} rooted at an odd
integer $u$ coprime to $q$ is built by, at each odd node $u$, adjoining
a child $w_a := (2^a u - 1)/q$ for every $a\ge 1$ such that
$2^a u \equiv 1 \pmod q$; such $w_a$ is automatically an odd integer
whenever it is an integer, since an odd numerator divided evenly by an
odd denominator yields an odd quotient. Writing $d:=\ord_q(2)$ for the
multiplicative order of $2$ modulo $q$, admissibility of $a$ for a
given parent $u$ depends only on $u\bmod q$: \emph{if} some
$a_0(u)\in\{1,\dots,d\}$ satisfies $2^{a_0(u)}u\equiv 1\pmod q$, it is
unique, and the admissible exponents are exactly $a\equiv a_0(u)\pmod
d$. Such $a_0(u)$ exists exactly when $u\bmod q$ lies in the subgroup
$\langle2\rangle\le(\mathbb{Z}/q\mathbb{Z})^\times$ generated by $2$;
nodes with $u\equiv 0\pmod q$, and, more generally, nodes whose residue
mod $q$ falls outside $\langle2\rangle$, are \emph{sterile} (no
admissible children). When $2$ is a primitive root mod $q$ (as for
$q=3,5,9$, all used numerically below), $\langle2\rangle$ is all of
$(\mathbb{Z}/q\mathbb{Z})^\times$ and $u\equiv 0\pmod q$ is the only
sterile class; for $q=3$ this recovers the familiar fact that
multiples of $3$ have trivial reverse subtrees. In general, however,
$d=\ord_q(2)$ can be a proper divisor of $\varphi(q)$: at $q=7$,
$d=3<\varphi(7)=6$ and $\langle2\rangle=\{1,2,4\}$, so $u\equiv
3,5,6\pmod7$ are sterile too, not merely $u\equiv 0$. This does not
affect the pressure identity of Theorem~\ref{thm:pressure} below (whose
proof, via Lemma~\ref{lem:fibre-bijection}, works in the reverse
direction --- from an unconstrained exponent sequence to its unique
root --- and never invokes existence of $a_0(u)$ for every $u$), nor
the numerical enumeration of \S\ref{sec:kl-empirical}, whose
implementation already excludes exactly these classes; it is a
correction to this structural description only.

We count distinct integer labels rather than path occurrences. Define
\[
 N_u(x):=\#\{n\le x:n\text{ odd and }T_q^k(n)=u
                    \text{ for some }k\ge0\}.
\]
This convention removes repetitions caused by cycles. In particular,
$N_u(x)\le\lceil x/2\rceil$.

On this tree we study the linear population functional $G$ satisfying
the fixed-point recursion
\begin{equation}\label{eq:G-recursion}
  G(u) = \sum_{a \equiv a_0(u) \ (d)} q\cdot 2^{-a}\, G(w_a),
\end{equation}
which for $q=3$ is exactly the martingale-type recursion underlying the
$G(v)/\mu$ line of investigation summarized in \S\ref{sec:endogeny}
below, and which is the direct reverse-tree analogue of
\eqref{eq:syracuse-def}.

\section{The closed-form pressure equation}
\label{sec:pressure}

\subsection{The pressure operator: an exact annealed identity}

A first, natural attempt to organize the recursion \eqref{eq:G-recursion}
into a transfer operator is to track the parent's residue class modulo
$q$ alone. This attempt fails, and it is worth recording exactly why,
since an earlier draft of this paper asserted a version of it that does
not survive scrutiny.

\begin{example}[The naive mod-$q$ automaton is not well defined]
\label{ex:naive-fails}
Take $q=3$ (so $d=\ord_3(2)=2$) and $a=2$, admissible for every parent
$u\equiv1\pmod3$. Two parents $u=1$ and $u=7$, both $\equiv1\pmod3$,
give $w_2(1)=(4-1)/3=1$ and $w_2(7)=(28-1)/3=9$, with $w_2(1)\equiv
1\pmod3$ but $w_2(7)\equiv0\pmod3$: the child's residue class depends
on $u\bmod9$, not merely on $u\bmod3$. In general, writing
$u=u_0+qk$ with $u_0=u\bmod q$, one finds $w_a\bmod q =
[(2^au_0-1)/q\bmod q] + [2^ak\bmod q]$: an explicit dependence on
$k=\lfloor u/q\rfloor\bmod q$. No finite modulus $q^k$ repairs this:
the same computation shows $w_a\bmod q^k$ depends on $u\bmod q^{k+1}$,
one digit further than what mod-$q^k$ knowledge of $u$ supplies. There
is therefore no well-defined finite-state transfer operator tracking a
single generation of the recursion on residue classes of any fixed
finite modulus.
\end{example}

What \emph{is} well defined, and is the correct replacement for the
naive picture of Example~\ref{ex:naive-fails}, is a bijection between
an entire admissible fibre at one level of the digit tower and the
full residue ring one level below.

\begin{lemma}[Fibre bijection]\label{lem:fibre-bijection}
Fix $a\ge1$ and $j\ge1$. The map $v\mapsto w_a(v):=(2^av-1)/q$ is a
bijection from $\{v\bmod q^j : 2^av\equiv1\!\!\pmod q\}$ (an arithmetic
progression of $q^{j-1}$ residues mod $q^j$) onto
$\mathbb{Z}/q^{j-1}\mathbb{Z}$.
\end{lemma}
\begin{proof}
Well defined: $v\bmod q^j$ determines $2^av-1\bmod q^j$, hence
$w_a(v)\bmod q^{j-1}$. Injective: if $v\not\equiv v'\pmod{q^j}$ within
the progression, write $v-v'=qt$ with $t\not\equiv0\pmod{q^{j-1}}$;
then $w_a(v)-w_a(v')=2^at\not\equiv0\pmod{q^{j-1}}$ since $2$ is a unit
mod $q^{j-1}$. Surjectivity follows since the map is an injection
between two finite sets of equal size $q^{j-1}$.
\end{proof}

Iterating Lemma~\ref{lem:fibre-bijection} along a depth-$k$ path
$(a_1,\dots,a_k)$ recovers, in truncated form, exactly the self-similar
identity \eqref{eq:X-selfsimilar} of \S\ref{sec:context} below: every
sequence $(a_1,\dots,a_k)\in(\mathbb{Z}_{\ge1})^k$, with no
admissibility restriction on the individual $a_i$, is admissible for
\emph{exactly one} root residue $u_0\bmod q^k$, namely
$u_0\equiv\sum_{i=1}^k q^{i-1}2^{-S_i}\pmod{q^k}$, $S_i:=a_1+\cdots+a_i$.
Define the depth-$k$ partition function rooted at $u_0$,
\[
  Z_k(\alpha;u_0) := \sum_{\substack{(a_1,\dots,a_k)\text{ admissible}\\ \text{from } u_0}} \ \prod_{i=1}^k \bigl(q\cdot2^{-a_i}\bigr)^\alpha .
\]

\begin{theorem}[Exact annealed pressure identity]\label{thm:pressure}
For every odd $q\ge3$, every $k\ge1$, and every $\alpha>0$,
\begin{equation}\label{eq:annealed-identity}
  \sum_{u_0\in\mathbb{Z}/q^k\mathbb{Z}} Z_k(\alpha;u_0) \;=\; \Bigl(\frac{q^\alpha}{2^\alpha-1}\Bigr)^{k},
\end{equation}
so the \emph{average} of $Z_k(\alpha;\cdot)$ over all $q^k$ root
residues equals exactly $\bigl(q^{\alpha-1}/(2^\alpha-1)\bigr)^k$; the
associated annealed pressure is
\begin{equation}\label{eq:pressure-closed-form}
  \rho_{\mathrm{ann}}(\alpha) = \frac{q^{\alpha-1}}{2^\alpha-1},
\end{equation}
giving the pressure equation at $\rho_{\mathrm{ann}}=1$:
\begin{equation}\label{eq:pressure-eq}
  q^{\alpha-1} = 2^{\alpha}-1.
\end{equation}
\end{theorem}
\begin{proof}
By Lemma~\ref{lem:fibre-bijection} (iterated), the admissible depth-$k$
sequences rooted at the $q^k$ distinct residues $u_0$ partition the
full, unconstrained set $(\mathbb{Z}_{\ge1})^k$ with no overlap and no
omission. Summing $\prod_i(q2^{-a_i})^\alpha$ over \emph{all} sequences
in $(\mathbb{Z}_{\ge1})^k$ therefore factors as
$\bigl(\sum_{a=1}^\infty (q2^{-a})^\alpha\bigr)^k =
\bigl(q^\alpha/(2^\alpha-1)\bigr)^k$, which is exactly the left-hand
side of \eqref{eq:annealed-identity} summed over the partition.
\end{proof}

We verified \eqref{eq:annealed-identity} independently by direct
enumeration for $q\in\{3,5,7,9\}$, $k\in\{1,2,3\}$, and
$\alpha\in\{0.26,0.5,1,1.37,2\}$, matching to the truncation error of
the (rapidly convergent) tail of the sum over $a$.

\begin{remark}[Mass conservation as a corollary]
The right-hand side of \eqref{eq:annealed-identity} evaluates to
exactly $1^k=1$ at $\alpha=1$ for every $q$ (mass conservation of the
branching process). Thus $\alpha=1$ is always a root of
\eqref{eq:pressure-eq}.
\end{remark}

\begin{theorem}[$q$-adic density martingale]\label{thm:qadic-martingale}
Let $\alpha$ be either positive root of \eqref{eq:pressure-eq}, and let
$U$ be Haar-uniform on $\mathbb Z_q$.
Then
\[
 M_k^{(\alpha)}(U):=Z_k(\alpha;U\bmod q^k)
\]
is a nonnegative martingale of mean one for the filtration generated by
$U\bmod q^k$. Consequently $M_k$ converges almost surely to a finite
limit $W_{q,\alpha}$. If $\mu_\alpha$ is the projective law of the tilted
Syracuse sums and
$\mu_\alpha=f_{q,\alpha}\,m_q+\mu_{\mathrm s}$ is its Lebesgue decomposition
with respect to Haar measure $m_q$, then
\[
 W_{q,\alpha}=f_{q,\alpha}(U)\qquad m_q\text{-almost surely}.
\]
In particular, $W_{q,\alpha}$ is nonzero with positive probability
exactly when $\mu_\alpha$ has a nonzero absolutely continuous component,
and $\mathbb E W_{q,\alpha}=1$ exactly when $\mu_\alpha\ll m_q$.
\end{theorem}

\begin{proof}
At a pressure root, the numbers
\[
 p_\alpha(a):=q^{\alpha-1}2^{-\alpha a},\qquad a\ge1,
\]
sum to one. For independent digits with this law, let
\[
 F_k=\sum_{i=1}^k q^{i-1}2^{-S_i}\pmod {q^k},
 \qquad S_i=a_1+\cdots+a_i.
\]
The fibre identity gives
\[
 Z_k(\alpha;u)=q^k\Pr(F_k=u).
\]
Also $F_{k+1}\equiv F_k\pmod {q^k}$, since their difference is
$q^k2^{-S_{k+1}}$. Thus the laws of $F_k$ form a projectively
consistent family. Conditional on $U\equiv u\pmod {q^k}$,
\[
 \mathbb E[M_{k+1}\mid U\bmod q^k=u]
 =\frac1q\sum_{v\equiv u\ (q^k)}q^{k+1}\Pr(F_{k+1}=v)
 =q^k\Pr(F_k=u)=M_k(u).
\]
The nonnegative martingale convergence theorem gives the limit. The
functions $M_k$ are the Radon--Nikodym derivatives of the restrictions
of $\mu_\alpha$ to the cylinder sigma-algebras relative to the
corresponding restrictions of $m_q$. Martingale differentiation of the
Lebesgue decomposition identifies their almost-sure limit with
$d(\mu_\alpha)_{\mathrm{ac}}/dm_q$. Its integral is the total mass of
the absolutely continuous part, which proves the final assertions.
\end{proof}

\begin{theorem}[$L^p$ collision criterion]\label{thm:lp-collision}
Let $\alpha$ be a pressure root, let $p>1$, let $\mu_{\alpha,k}$ be the
law of $F_k$ modulo $q^k$, and
let $M_k^{(\alpha)}=q^k\mu_{\alpha,k}$ be its density relative to
uniform measure on $\mathbb Z/q^k\mathbb Z$. Then
\begin{equation}\label{eq:lp-collision}
 \lVert M_k^{(\alpha)}\rVert_p^p
 =q^{k(p-1)}\sum_{u\bmod q^k}\mu_{\alpha,k}(u)^p.
\end{equation}
The projective measure $\mu_\alpha$ has a density in $L^p(\mathbb Z_q)$
if and only if the right-hand side is bounded uniformly in $k$. If $p$
is an integer, the sum in \eqref{eq:lp-collision} is the probability
that $p$ independent tilted Syracuse sums agree modulo $q^k$.
\end{theorem}

\begin{proof}
Uniform measure assigns mass $q^{-k}$ to each residue, so
\[
 \lVert M_k^{(\alpha)}\rVert_p^p
 =q^{-k}\sum_u\bigl(q^k\mu_{\alpha,k}(u)\bigr)^p,
\]
which is \eqref{eq:lp-collision}. For integer $p$, expanding the
probability of equality of $p$ independent copies gives the collision
interpretation.

If the quantities in \eqref{eq:lp-collision} are uniformly bounded,
the martingale $M_k^{(\alpha)}$ converges in $L^p$. Its limit has integral
one and is the density of the projective measure on every cylinder, so
$\mu_\alpha$ is absolutely continuous with an $L^p$ density. Conversely,
if $d\mu_\alpha=f\,dm_q$ with $f\in L^p$, then
$M_k^{(\alpha)}=\mathbb E[f\mid U\bmod q^k]$. Conditional expectation is
an $L^p$ contraction, which gives the uniform bound.
\end{proof}

For $q=3$ and $\alpha=1$, Theorem~\ref{thm:lp-collision} at $p=2$
gives exactly the collision statistic $K_k$ used in
\S\ref{sec:lemmas}. Thus the former $L^2$ sufficient condition is one
member of a complete family of weighted collision criteria (finite-level
spectrum computed in \texttt{collatz-endogeny/\allowbreak
sec3-pressure-equation}). These are
averaged concentration conditions. They do not control the smallest
cylinder mass required by the worst-residue problem $\beta=1$.

\begin{theorem}[Singularity at the frozen root]
\label{thm:frozen-singular}
Let $\alpha=\alpha_+(q)$. Then
\[
 M_k^{(\alpha)}(U)\longrightarrow0
 \qquad m_q\text{-almost surely},
\]
and the tilted Syracuse measure $\mu_\alpha$ is singular with respect
to Haar measure.
\end{theorem}

\begin{proof}
Write $r=2^{-\alpha}$. The digit law is
$p_\alpha(a)=(1-r)r^{a-1}$, so its entropy is
\[
 H(p_\alpha)=-\log(1-r)-\frac{r}{1-r}\log r.
\]
At a pressure root, direct differentiation gives
\[
 H(p_\alpha)-\log q
 =-\alpha P'(\alpha)=s(\alpha).
\]
The larger root is frozen, hence $H(p_\alpha)<\log q$. For $0<t<1$,
subadditivity of $x^t$ and the representation
$M_k^{(\alpha)}(u)=q^k\sum_{F_k=u}\prod_i p_\alpha(a_i)$ give
\[
 \mathbb E_{m_q}\bigl[(M_k^{(\alpha)})^t\bigr]
 \le\left(q^{t-1}\sum_{a\ge1}p_\alpha(a)^t\right)^k.
\]
The logarithmic derivative of the factor in parentheses at $t=1$ is
$\log q-H(p_\alpha)>0$. It is therefore less than one for some
$t<1$ sufficiently close to one. The fractional moment tends to zero,
so the almost-sure martingale limit is zero. Theorem~\ref{thm:qadic-martingale}
then identifies the absolutely continuous component as zero.
\end{proof}

\subsection{Quenched vs.\ annealed: a freezing transition}
\label{sec:freezing}

Theorem~\ref{thm:pressure} is a statement about the \emph{average} of
$Z_k(\alpha;\cdot)$ over root residues; it is not, by itself, a
statement about $Z_k(\alpha;u_0)^{1/k}$ for a single, fixed root
$u_0$ --- which is what actually governs the population martingale
$G(v)$ for one specific integer $v$. Whether the two coincide is
exactly the quenched-versus-annealed question familiar from directed
polymers in a random environment and from Derrida's random energy
model, governed by a standard freezing criterion. Writing
$P(\alpha):=\log\rho_{\mathrm{ann}}(\alpha)$ for the annealed
log-pressure and $s(\alpha):=P(\alpha)-\alpha P'(\alpha)$ for the
associated entropy, quenched and annealed rates coincide where
$s(\alpha)>0$ (an \emph{unfrozen} phase) and diverge where
$s(\alpha)<0$ (a \emph{frozen} phase, in which the quenched rate is
dominated by rare, low-entropy paths rather than by the bulk of the
ensemble).

\begin{proposition}[The larger root is always frozen]\label{prop:always-frozen}
For every odd $q\ge3$, writing $\alpha_-<\alpha_+$ for the two positive
roots of $\rho_{\mathrm{ann}}(\alpha)=1$: $s(\alpha_-)>0$ (unfrozen) and
$s(\alpha_+)<0$ (frozen), unconditionally.
\end{proposition}
\begin{proof}
At any root, $P(\alpha)=0$, so $s(\alpha)=-\alpha P'(\alpha)$. Since
$\rho_{\mathrm{ann}}=e^P$ is log-convex (\S\ref{sec:endogeny-addendum}),
$P$ itself is strictly convex; with exactly two zeros
$\alpha_-<\alpha_+$, necessarily $P'(\alpha_-)<0<P'(\alpha_+)$
(the function must be decreasing through the first zero and increasing
through the second, on either side of its unique minimum). Hence
$s(\alpha_-)=-\alpha_-P'(\alpha_-)>0$ and
$s(\alpha_+)=-\alpha_+P'(\alpha_+)<0$.
\end{proof}

For $q\ge5$, where $\alpha_+=1$ always, this has an exact closed form:
at $\alpha=1$, $2^\alpha-1=1$ so $\log(2^\alpha-1)=0$, and
$s(1)=2\log2-\log q=\log(4/q)$, which is negative exactly when $q>4$.
So the freezing of the trivial root for every odd $q\ge5$ is not a
numerical coincidence but the same sign condition as the mean-drift
remark above ($\log q-2\log2>0$): the drift that makes orbits expand on
average is, up to a minus sign, exactly the entropy deficit that
freezes $\alpha_+=1$.

We computed the freezing threshold $\alpha_c(q)$ (the root of
$s(\alpha)=0$, which by Proposition~\ref{prop:always-frozen} always
lies strictly between $\alpha_-(q)$ and $\alpha_+(q)$) directly:
\begin{center}
\begin{tabular}{c|c}
$q$ & $\alpha_c(q)$ \\ \hline
$3$ & $1.355$ \\
$5$ & $0.794$ \\
$7$ & $0.564$ \\
$9$ & $0.437$
\end{tabular}
\end{center}
confirming $\alpha_-(q)<\alpha_c(q)<\alpha_+(q)$ in every case, as
Proposition~\ref{prop:always-frozen} requires. We also verified
numerically, with an independent memoised recursion computing
$Z_k(\alpha;u_0)$ exactly at a fixed (not averaged) root, that the
freezing is not merely a formal entropy computation but a genuine
divergence of quenched from annealed growth: at $q=3$, $u_0=1$ --- a
genuine period-one cycle of the residue dynamics, since $w_2(1)=1$
exactly --- $Z_k(2;1)^{1/k}$ shows an oscillation with geometric-mean
ratio well below the annealed rate $\rho_{\mathrm{ann}}(2)=1$; at
$q=5$, $Z_k(1;u_0)^{1/k}$ at the roots $u_0\in\{1,2,3,4\}$ likewise
fails to track $\rho_{\mathrm{ann}}(1)=1$. We flag, rather than
obscure, a limitation of this specific check: every one of these
roots lies on or feeds directly into the smallest-admissible-branch
map's own trivial cycle ($1\leftrightarrow3$ for $q=5$, exactly as
$u_0=1$ does for $q=3$), so the divergence observed here has a
deterministic-recurrence explanation available in addition to the
entropy mechanism of Proposition~\ref{prop:always-frozen}, and does
not by itself sample the generic (non-periodic) environment that the
quenched statement is really about. We therefore report this as an
illustration of freezing at a special, easily-computed class of
roots, not as a confirmation at a generic one; \S\ref{sec:kl-empirical}
below and the statistically calibrated measurement of \S\ref{sec:kl}
--- both averaging over many roots sampled well away from any known
short cycle --- carry the actual generic-environment evidence.

The pressure identity transfers to the counting exponent of the i.i.d.
model at the unfrozen root $\alpha_-(q)$. It does not by itself prove a
tail law. Theorem~\ref{thm:iid-tail} obtains the i.i.d. tail by implicit
renewal, while Conjecture~\ref{conj:tail-index} states the separate
claim for a coherent Haar-random $q$-adic root environment.

\subsection{The structural transition at $q=5$}

Besides the trivial root $\alpha=1$, equation~\eqref{eq:pressure-eq}
has exactly one further positive root $\alpha^\ast(q)$ (by strict
log-convexity of $\alpha\mapsto q^{\alpha-1}/(2^\alpha-1)$ on
$\alpha>0$, discussed in \S\ref{sec:endogeny-addendum} below, the
equation $\rho(M_q(\alpha))=1$ has at most two positive roots). We
solved for $\alpha^\ast(q)$ independently by bisection:

\begin{center}
\begin{tabular}{c|c|c}
$q$ & second root $\alpha^\ast(q)$ & interpretation \\ \hline
$3$ & $2.000000$ & exact \\
$5$ & $0.650919$ & \\
$7$ & $0.373501$ & \\
$9$ & $0.258108$ &
\end{tabular}
\end{center}

At $q=3$ the second root is exactly $\alpha^\ast=2$, matching the
universal tail exponent for the classical Collatz reverse tree
established by independent methods (Hill estimator on real-tree
samples, and extreme value theory on block maxima) in earlier stages
of this project. For $q\ge 5$, however, the second root falls
\emph{below} $1$ rather than above it: a genuine structural transition,
not a numerically small perturbation. Writing $\alpha_-<\alpha_+$ for
the two roots, we have $\{\alpha_-,\alpha_+\}=\{1,2\}$ at $q=3$ but
$\{\alpha_-,\alpha_+\}=\{\alpha^\ast(q),1\}$ for $q\ge5$: the trivial
root becomes the \emph{larger} of the two.

\begin{theorem}[Structural transition: exponent, i.i.d.\ model]\label{thm:transition-model}
Let $N(x)$ be the population counting function of the i.i.d.\
branching-random-walk model whose annealed pressure is exactly
$\rho_{\mathrm{ann}}(\alpha)=q^{\alpha-1}/(2^\alpha-1)$
(Theorem~\ref{thm:pressure}). Since $\alpha_-(q)$ is always unfrozen
(Proposition~\ref{prop:always-frozen}), $\log N(x)/\log x\to\alpha_-(q)$
almost surely, for every odd $q\ge3$ \cite{Biggins1992}: exponent $1$
(linear growth) at $q=3$; exponent $\alpha_-(q)<1$ for every $q\ge5$.
\end{theorem}

\begin{proposition}[Finer asymptotic, conditional]\label{prop:transition-fine}
Conditional on a renewal-type refinement, from the counting-function
level (not merely the exponent) to the sharper asymptotic
$N(x)\sim W\cdot x^{\alpha_-(q)}$, holding almost surely and uniformly,
\emph{including} at $q=3$, for an almost surely positive, non-degenerate
random scale factor $W$ (not a deterministic constant at any $q$, since
the underlying reproduction law has genuine variance). A directed
literature search found the neighboring machinery but not this exact
statement ready-made: the classical Crump--Mode--Jagers renewal theorem
\cite{Nerman1981} is stated for one-sided ($[0,\infty)$) reproduction
point processes, while our branching displacements include negative
values (e.g.\ $2/3<1$ at $q=3$, since a child can be numerically
smaller than its parent); the branching-random-walk martingale theory
that is natively two-sided \cite{Biggins1992} gives convergence
generation-by-generation, not a renewal law summed over all generations;
and the implicit renewal theory for weighted trees with general
(two-sided) weights \cite{Goldie1991,JelenkovicOlveraCravioto2012,
AlsmeyerBigginsMeiners2012} resolves the \emph{tail} of the fixed-point
equation for $W$ itself, not the cumulative counting law for $N(x)$ ---
a related but distinct question. We expect the present statement to
follow from combining this machinery (spine decomposition and
exponential tilting for a two-sided one-dimensional random walk,
exactly as already used for the tail-of-$W$ results just cited) but we
have not carried out that adaptation, nor located it already written in
this exact form; a very recent preprint \cite{VillemonaisZalduendo2025}
extends branching-process convergence laws to general type spaces
without geometric rescaling and may already contain it, but we have not
verified this. The non-lattice hypothesis that any such extension of
\cite{Nerman1981} would need holds here, and it is precisely the
irrationality of $\log q/\log2$ that secures it. Although each
single-step displacement $-\log q + a\log2$ ($a\in\mathbb{Z}_{\ge1}$)
lies in one coset of the lattice $\log(2)\,\mathbb{Z}$, containment in
a coset is not arithmeticity: the smallest closed subgroup of
$\mathbb{R}$ generated by these displacements is
$\log(q)\,\mathbb{Z}+\log(2)\,\mathbb{Z}$, which is dense in
$\mathbb{R}$ for every odd $q\ge3$ because $\log q/\log2\notin\mathbb{Q}$.
Equivalently, the characteristic function of the displacement equals
$1$ only at the origin, so the renewal/CMJ asymptotic carries no
log-periodic (period $\log2$) modulation. Were $\log q/\log2$ rational
the branch positions would collapse onto a genuine lattice and a
periodic correction would appear; the irrationality is therefore the
operative fact, not an irrelevant one.
\end{proposition}

\begin{conjecture}[Structural transition: exponent, arithmetic tree]\label{conj:transition-arithmetic}
For every admissible root $u$, the genuine reverse tree of $T_q$
satisfies
\[
 N_u(x)=x^{\alpha_-(q)+o(1)}.
\]
Thus the exponent is $1$ for $q=3$ and is strictly below $1$ for every
odd $q\ge5$. At $q=3$ this is the Growth Exponent Conjecture of
Kontorovich--Lagarias \cite{KontorovichLagarias2009} and
Applegate--Lagarias
\cite{ApplegateLagarias1995I,ApplegateLagarias1995II} for the classical
Collatz reverse tree, namely $\pi_a(x)=x^{1+o(1)}$, open
since 1995, with the best unconditional partial result remaining the
one-sided bound $\pi_a(x)\ge x^{0.84}$ of Krasikov--Lagarias
\cite{KrasikovLagarias2003}. Our pressure identity reproduces the
exponent-$1$ statement of this conjecture as the $\alpha_-=1$ case of
a single closed-form equation, together with its $q\ge5$ analogue,
rather than resolving it.
\end{conjecture}

We state Theorem~\ref{thm:transition-model},
Proposition~\ref{prop:transition-fine}, and
Conjecture~\ref{conj:transition-arithmetic} separately, rather than as
a single theorem about the arithmetic tree, precisely because the step
from one to the other is exactly the kind of transfer that caused the
error corrected in this section; Remark~\ref{rem:transfer-basis}
below makes precise what does and does not rest on proof. Direct
empirical support for Conjecture~\ref{conj:transition-arithmetic} is
given in \S\ref{sec:kl-empirical} below and by the statistically
calibrated measurement of \S\ref{sec:kl}.

\begin{remark}[What ``transfers'' rests on]\label{rem:transfer-basis}
We are precise about the status of this transfer, since it is exactly
the kind of step that caused the error corrected in this section. For
the genuine i.i.d.\ branching-random-walk model that the annealed
identity of Theorem~\ref{thm:pressure} describes, quenched-equals-annealed
growth throughout the unfrozen region is a rigorous, classical result
\cite{Biggins1992}. For our arithmetic recursion, no analogous theorem
is known to us: the gap is the same one already identified around
Theorem~\ref{thm:barrier} (Remark~\ref{rem:barrier-reading}:
cancellation of centered fine-frequency correlations after path
averaging, here in a quantitative large-deviations form), and
the closest genuine arithmetic result in
the literature is a one-sided density bound of Applegate--Lagarias
\cite{ApplegateLagarias1995I,ApplegateLagarias1995II}, later sharpened
by Krasikov--Lagarias \cite{KrasikovLagarias2003} (already cited in
\S\ref{sec:triangulation}), not a matching two-sided quenched
statement. We therefore treat the transfer, at every root we invoke it
for, as supported by (i) the classical stochastic-model theorem, (ii)
direct numerical confirmation, and not by a proof for the arithmetic
tree itself --- naming the gap precisely rather than letting the
stochastic-model theorem's rigor bleed by association into the
arithmetic claim.
\end{remark}

\begin{conjecture}[Density and tail of the tilted Syracuse measure]
\label{conj:tail-index}
Let $U$ be Haar-uniform on $\mathbb Z_q$, and read
$Z_k(\alpha_-(q);U)$ through its residue modulo $q^k$. The sequence
has the limit $W_q:=W_{q,\alpha_-(q)}$ of
Theorem~\ref{thm:qadic-martingale}. The tilted
measure is absolutely continuous, its density $W_q$ is non-degenerate,
and
\[
 \Pr(W_q>x)=x^{-\alpha_+(q)/\alpha_-(q)+o(1)}.
\]
Here $\alpha_+(q)$ is the larger root of \eqref{eq:pressure-eq}
($\alpha_+(3)=2$; $\alpha_+(q)=1$ for every $q\ge5$).
\end{conjecture}

\begin{theorem}[Tail index in the i.i.d. model]\label{thm:iid-tail}
For the additive martingale at parameter $\alpha_-(q)$ in the matching
i.i.d. branching model,
\[
 \Pr(W_\infty>x)\sim C_q x^{-\alpha_+(q)/\alpha_-(q)}
 \qquad (x\to\infty)
\]
for a constant $C_q>0$.
\end{theorem}

\begin{proof}
Apply the implicit renewal theorem for generalized multiplicative
cascades \cite{Liu2000}. With $\theta=\alpha_-(q)$, its logarithmic
moment function is
\[
 \psi(s)=\log\rho_{\mathrm{ann}}(\theta s).
\]
It vanishes at $s=1$ and next at
$\kappa=\alpha_+(q)/\alpha_-(q)$, while $\psi'(1)<0$ because the smaller
root is unfrozen. The geometric offspring ladders give all moments
needed in a neighborhood of $\kappa$. The closed subgroup generated by
the logarithmic weights is generated by $\theta\log2$ and
$\theta\log q$; it is dense because $\log q/\log2$ is irrational.
Thus the renewal law is non-arithmetic and has a constant, positive
tail coefficient.
\end{proof}

The four hypotheses above were checked numerically for
$q=3,5,7,9,15$ in \texttt{collatz-endogeny/\allowbreak
sec3-pressure-equation}.

This is the classical smoothing-transform result for the stochastic
model. Its transfer to the coherent $q$-adic root environment is
Conjecture~\ref{conj:tail-index}, since the arithmetic subtrees are not
independent.\footnote{A related
structural prediction of the multitype model,
confirmed empirically for $q=5$: conditioning on the root's residue
class $i=u_0\bmod q$ (with $a_0(i)$ the smallest admissible branch
exponent for class $i$), the scale factor satisfies
$W_i\stackrel{d}{=}2^{-a_0(i)\theta}\cdot W^*$ for a common $W^*$,
i.e.\ residue class only rescales the candidate limit, it does not
change its shape.
Verified on $5000$ sampled roots across four headroom levels, with
quantile ratios matching this prediction to within $2$--$9\%$; exact
only in the idealized i.i.d.\ model, as in Remark~\ref{rem:transfer-basis}.
The same relation survives unchanged when $2$ is not a primitive root
mod $q$ (the extra sterile classes of \S\ref{sec:setup}): restricted
to the surviving types $\langle2\rangle$, the multitype pressure
matrix is rank $1$, and its Perron eigenvalue collapses, by a
$d$-independent cancellation ($a_0$ a bijection
$\langle2\rangle\to\{1,\ldots,d\}$), to exactly $q^{s-1}/(2^s-1)$ ---
the same pressure equation --- so both the exponent and the
scale-family direction survive intact; verified for $q=7,15$ to
$1$--$4\%$. (Whether $2$ is a primitive root mod $q$ is, for prime
$q$, governed by Artin's primitive root conjecture, open in general
but proved conditionally on GRH by \cite{Hooley1967}.)}
the tail-index equation $\mathbb{E}[\sum_i A_i^s]=1$ becomes,
for our weights, $q^{\alpha_-s-1}=2^{\alpha_- s}-1$, solved exactly by
$s=\alpha_+/\alpha_-$. This is further reason to take the conjecture
seriously as a prediction, even though, per
Remark~\ref{rem:transfer-basis}, the theory only rigorously covers the
stochastic model, not our arithmetic recursion.

We state this as a conjecture separate from both
Theorem~\ref{thm:transition-model} and
Conjecture~\ref{conj:transition-arithmetic} above, precisely because
Proposition~\ref{prop:always-frozen} shows $\alpha_+(q)$ is
\emph{always} frozen, so the annealed pressure identity alone does not
establish the tail. Theorem~\ref{thm:iid-tail} obtains it by implicit
renewal at the unfrozen martingale normalization. Its arithmetic
counterpart still requires evidence external to
Theorem~\ref{thm:pressure}. At $q=3$ that evidence is real and
threefold: the Hill estimator on real-tree samples and extreme-value
theory on block maxima (both from earlier stages of this project) give
$\alpha\approx2$ independently of the pressure argument, and the
counting-exponent confirmation below adds a third, related line of
support. For $q\ge5$ the corresponding evidence is markedly weaker: an
early Hill-estimator measurement of finite-depth root scale factors for
$q=5$ (600
roots, agreement to two decimal places with the predicted
$1.536\ldots$) was later shown, over the course of this project, to be
statistically non-confirmatory at that sample size (the estimator's
true standard error is $\approx0.45$, an order of magnitude larger
than the agreement being claimed). We pursued this further along two
independent routes, and report both honestly rather than stopping at
whichever gave the more favorable number.

First, a substantially larger statistical test (5000 roots, four
headroom levels up to $10^8$; \texttt{collatz-endogeny/\allowbreak sec3-pressure-equation})
applied a battery of four estimators (a rank-size regression with the
Gabaix--Ibragimov finite-sample correction; a bias-corrected Hill
estimator following Huisman et al.; a generalized-Pareto threshold-stability
scan; and a Clauset--Shalizi--Newman fit with a Vuong test against a
truncated log-normal alternative). The outcome is mixed rather than
confirmatory: the bias-corrected Hill estimator is remarkably stable
across headroom levels and lands close to the predicted $1.536$, but
the threshold-stability scan shows no clean plateau near the predicted
shape parameter, and the Vuong test favors the log-normal alternative
over the power law, at conventional significance, in three of the four
headroom levels tested. Reading the four estimators together by the
depth of tail they each summarize shows the apparent local tail index
rising smoothly with depth (from $\approx1.3$ at the widest window to
$\approx2.2$ at the narrowest) rather than settling near any single
value --- consistent with a slowly-converging pre-asymptotic regime in
which no estimator at this sample size can be treated as decisive.

Second, we tested the conjecture by an exact, non-statistical route:
the same depth-$k$ recursion $Z_k(\theta;u)$ used to verify the
frozen-phase growth law above (with $\theta=\alpha_-(5)$) was evaluated
over the \emph{entire} population of residues $u\bmod 5^k$ (not a
sample), for $k$ up to $11$ ($5^{11}\approx4.9\times10^7$ residues; a
sanity check --- that the population mean of $Z_k(\theta;u)$ equals
exactly $1$ for every $k$, forced by Theorem~\ref{thm:pressure} at
$\alpha=\theta$ --- passed to floating-point precision, confirming the
implementation; \texttt{collatz-endogeny/\allowbreak sec3-pressure-equation}).
If $W_q=\lim_k Z_k(\theta;U)$ has tail index
$\iota$, the population moment $M_k(p)=\mathrm{mean}_u\,Z_k(\theta;u)^p$
should saturate in $k$ for $p<\iota$ and diverge for $p>\iota$.
The data show saturating behavior throughout $p\le1.6$ and diverging
behavior for $p\ge1.7$, which would place $\iota$ higher than the
predicted $1.536$ if taken at face value --- but the ratio of
successive increments has not yet stabilized for $p\le1.6$ at $k=11$
(unlike $p\ge1.7$, where it has), the classic signature of a system
still relaxing toward its asymptotic exponent rather than one that has
reached it. Given the previously observed slow convergence rate for
$q=5$ ($k^{-0.222}$ decay in the moment test --- not an effect of the
linear transfer operator's spectrum, which we later verified to be
exactly $\{\Lambda,0\}$ at every truncation level, i.e. a perfect
spectral gap with no isolated subdominant eigenvalue; the transient's
origin lies in a separate, nonlinear critical layer and remains open),
the transient decays by only about $7\%$ between
$k=8$ and $k=11$; halving it would require $k\approx250$, far beyond
what exhaustive enumeration over $5^k$ residues can reach. We therefore
cannot distinguish, from this data, between $\theta=1.536$ lying just
below a nearby true index and lying further below one --- the test is
inconclusive, not disconfirming, and enumeration alone will not resolve
it.

Taken together, neither route confirms Conjecture~\ref{conj:tail-index}
for $q\ge5$ and neither refutes it; both are more informative than the
original under-powered measurement, and both point to the same honest
conclusion: this is a harder statistical question than the $q=3$ case,
which has stronger finite-level support from the independent Hill/EVT
measurements noted in \S\ref{sec:pressure}. We have since pursued exactly
the two follow-ups this pointed to. A pre-registered log-periodic fit
(motivated by the branch weights being powers of $2$) was tested against
a period derived analytically, not fitted: the relevant renewal-theoretic
dichotomy is non-arithmetic (the shift $\log_2 5$ is irrational), so no
asymptotic log-periodicity is predicted, and a periodogram test found no
power at the predicted frequencies. An analytic handle on the linear
transfer operator's spectrum, once obtained, likewise turned out not to
bear on this question, since the transient lives in a separate nonlinear
layer (the operator has a provably perfect spectral gap). The other
follow-up did move the needle: a pre-registered comparison on a
$20\times$ larger sample ($10^5$ roots) produced, for the first time,
a clean GPD threshold-stability plateau at the predicted value, a
tightly concentrated Hill estimate, and a Vuong test that no longer
favors the lognormal alternative -- all confirmed against two synthetic
calibration checks rather than taken at face value. This is now the
strongest evidence to date in favor of Conjecture~\ref{conj:tail-index}
for $q=5$, though not a proof of it: Vuong non-rejection is not
confirmation, and the martingale is provably still pre-asymptotic at
the headroom reached. We list the conjecture as an explicit open
problem in \S\ref{sec:conclusion}, updated to reflect this shift.

The mean drift $\log q - 2\log 2$ changes sign exactly at the
threshold $q=5$ separating the two cases of
Conjecture~\ref{conj:transition-arithmetic}: negative for $q=3$ (orbits contract on
average), positive for $q\ge 5$ (orbits expand on average). This is not
a coincidence alongside the freezing computation of
\S\ref{sec:freezing}: at $\alpha=1$ the entropy is exactly
$s(1)=\log(4/q)$, so the same sign that decides contraction versus
expansion of orbits also decides whether the trivial root $\alpha=1$
is frozen.

\subsection{Empirical confirmation on real reverse trees}
\label{sec:kl-empirical}

Independent depth-first enumeration of the true (not extended) reverse
trees for $q=5$ and $q=7$, following the exact admissibility rule of
\S\ref{sec:setup}, provides direct empirical support for
Conjecture~\ref{conj:transition-arithmetic}:
counting slopes per decade of $3$ independently sampled roots gave
$\approx 0.62$--$0.66$ for $q=5$ (predicted $0.650919$) and
$\approx 0.36$--$0.39$ for $q=7$ (predicted $0.373501$). A more careful,
statistically calibrated measurement for $q=5$ is the subject of
\S\ref{sec:kl}. As already noted above for $q=3$ (where $w_2(1)=1$ is a
genuine fixed point), the smallest-branch map $u\mapsto
w_{A_0(u\bmod q)}(u)$ has finite cycles for $q=5$ (e.g.\
$1\mapsto3\mapsto1$, since $w_4(1)=3$ and $w_1(3)=1$) and $q=7$; naive
enumeration would loop on nodes that reach one. Our implementation
excludes all such cycle members explicitly (precomputed for each $q$)
and deduplicates visited nodes, rather than relying only on sampling
roots far from them.

\begin{remark}[Bibliographic calibration]\label{rem:novelty-109}
A directed novelty check against Kontorovich--Lagarias
\cite{KontorovichLagarias2009}, Wirsching \cite{Wirsching1998}, and a
later rigorous forward-direction result excluding $q\ge5$ for the
analogous $p=2$ generalization \cite{GoncalvesGreenfeldMadrid2022}
established that neither the numerical value $\alpha^\ast(5)=0.650919$
nor the qualitative transition at $q\ge 5$ is new: Wirsching had
already conjectured the transition heuristically in 1998, and
Kontorovich--Lagarias had already computed the same quantity
(denoted $\eta_{5,BP}$ in their Theorem~8.10) via large-deviations
methods for a branching random walk, a route different from, but
numerically identical to, our pressure-equation derivation. Given
\S\ref{sec:freezing}, this difference in route may matter more than a
mere alternative derivation: a genuine large-deviations argument for a
branching random walk is, in the appropriate sense, a quenched
statement, whereas our pressure identity (Theorem~\ref{thm:pressure})
is annealed. That the two agree numerically at $\alpha_-(5)$ is
consistent with $\alpha_-(5)$ being unfrozen
(Proposition~\ref{prop:always-frozen}), where quenched and annealed
are expected to coincide; it does not by itself certify that our
annealed route would agree with a quenched one at a frozen root. What
survives as new here is the closed-form equation~\eqref{eq:pressure-eq}
itself, unifying the family for arbitrary $q$ in a single algebraic
identity, together with its rigorous annealed derivation and
independent numerical confirmation.
\end{remark}

\section{The endogeny barrier}
\label{sec:endogeny}

\subsection{Gauge freedom in the tree recursion}

The recursion \eqref{eq:G-recursion} is linear in $G$, and linearity
has an immediate structural consequence that is easy to overlook: if
$W$ denotes the ``arithmetic'' solution --- the one actually realized
by, and measurable with respect to, the $3$-adic digits of the tree
--- and $Y$ is \emph{any} bounded positive function invariant along
the tree (in particular, any function of an auxiliary coordinate
carried unchanged from parent to child), then $X := W\cdot Y$ solves
\emph{exactly the same recursion}.

\begin{proposition}[Gauge freedom]\label{prop:gauge}
Let $\mathbb{Z}_3^\ast\times\mathcal{Y}$ be the product space on which
the $3$-adic digit dynamics acts on the first factor exactly as on the
reverse tree, while passing an auxiliary coordinate $y\in\mathcal{Y}$
unchanged to every child. For bounded, non-constant $Y:\mathcal{Y}\to
(0,\infty)$, the function $X(r,y):=W(r)\cdot Y(y)$ satisfies:
\begin{enumerate}[label=(\alph*)]
\item the same pressure equation~\eqref{eq:pressure-eq} (identical
  weights, identical operator $M_q(\alpha)$, identical roots
  $\alpha=1,\alpha^\ast$);
\item the same tail index as $W$;
\item the same linear regression slope $b\approx 1$ against the
  limiting measure $\mu$, up to a fixed additional dispersion;
\item but $\Var[\log X \mid \text{first } m \text{ digits}] \to
  \Var[\log Y] > 0$ for \emph{every} $m$, including $m=\infty$: the
  conditional variance never collapses.
\end{enumerate}
\end{proposition}

Item (d) is the crux: no statistic computed from pressure, tail index,
or regression slope against $\mu$ distinguishes $W$ from $W\cdot Y$.
The technical name for this phenomenon is \emph{endogeny}
\cite{AldousBandyopadhyay2005}: collapse of the conditional variance is
equivalent to $X$ being measurable with respect to the $\sigma$-algebra
generated by the digits (``endogenous''); Proposition~\ref{prop:gauge}
exhibits a non-endogenous solution. We emphasize the logical status of
this construction: it proves that the recursion \emph{alone} does not
force exactness; it does not assert that a non-trivial $Y$ actually
occurs for the arithmetic object built from a single integer $v$ (there,
$Y$ would have to be manufactured from $v$ itself, not injected as an
independent auxiliary coordinate). These two statements must never be
conflated.

\subsection{A specific mechanism ruled out}

One natural candidate mechanism for a non-trivial gauge factor is a
persistent $2$-adic ``memory'' of the root surviving deep into the
tree. This is false:

\begin{lemma}[No $2$-adic memory]\label{lem:2adic}
Let $w$ be a child of $v$ at exponent $a$ in the reverse tree of the
classical ($q=3$) map, so $3w=2^a v - 1$. Then
\[
  w \equiv -3^{-1} \pmod{2^a},
\]
a universal constant depending only on $a$, independent of $v$.
Consequently, a descendant reached after a path with exponent sum $A$
has its lowest $A$ bits determined \emph{entirely by the path}, with
zero information about the root.
\end{lemma}

\begin{proof}
From $3w=2^av-1$ we get $3w\equiv -1\pmod{2^a}$, i.e.\ $w\equiv
-3^{-1}\pmod{2^a}$ since $3$ is a unit mod $2^a$. Induction along the
path, composing the relation at each generation, gives the claim.
\end{proof}

We verified Lemma~\ref{lem:2adic} numerically for $a=1,\dots,10$ against
more than $500$ values of $v$ per exponent, confirming the predicted
universal constant exactly in every case (e.g.\ $a=5\Rightarrow w\equiv
21\ (\mathrm{mod}\ 32)$ always; $a=10\Rightarrow w\equiv 341\
(\mathrm{mod}\ 1024)$ always). The reverse tree destroys all $2$-adic
information about the root at every generation; there is no persistent
$2$-adic coordinate available to correlate with the $3$-adic residue,
ruling out this specific mechanism for a gauge factor.

\subsection{What forces collapse, and what does not}

In the fully idealized probabilistic model --- i.i.d.\ branching with
independent subtree contents --- exactness \emph{is} forced, by two
classical mechanisms:

\begin{enumerate}[label=(\alph*)]
\item \textbf{Classification of smoothing-transform fixed points}
  \cite{DurrettLiggett1983,AlsmeyerBigginsMeiners2012}: non-endogenous
  solutions of the associated smoothing-transform fixed-point equation
  generically have a strictly smaller tail index than the endogenous
  one (in the regime relevant here, tail index $1$ with infinite mean),
  which is excluded by our own data: $\alpha^\ast=2$ has been measured
  by three independent methods (pressure equation, Hill estimator,
  extreme-value theory on block maxima).
\item \textbf{Choquet--Deny rigidity} for the Archimedean phase:
  bounded harmonic functions of a random walk on an abelian group are
  constant, which eliminates any Archimedean gauge coordinate under
  the full i.i.d.\ hypothesis.
\end{enumerate}

\begin{theorem}[The recursion does not force endogeny]\label{thm:barrier}
The smoothing-transform model associated with the exact weight
recursion of the reverse tree admits non-endogenous solutions (harmonic
gauge factors, Proposition~\ref{prop:gauge}; the $2$-adic obstruction is
excluded by Lemma~\ref{lem:2adic}) with a strictly positive
conditional-variance floor $\sigma_Y>0$ at every resolution,
statistically indistinguishable from the endogenous solution by
pressure, tail index, or any marginal we have tested. Consequently the
exact weight recursion together with the statistics of the $3$-adic
digits does not by itself force the residual gauge dispersion to
vanish.
\end{theorem}

\begin{remark}[Barrier reading, and epistemic status]\label{rem:barrier-reading}
Informally this delimits a class of arguments: we are aware of no
argument using only \textup{(i)} the exact weight recursion and
\textup{(ii)} statistics of the $3$-adic digits that does not factor
through this smoothing-transform model, and hence none that can
distinguish the endogenous from the non-endogenous solution. We do
\emph{not} formalize ``uses only (i) and (ii)'' as a precise class of
techniques (in contrast to relativization or natural-proofs barriers in
complexity theory), so this is a heuristic delineation of known
approaches, not a formal impossibility theorem. The $2$-adic exclusion
(Lemma~\ref{lem:2adic}) and the gauge construction
(Proposition~\ref{prop:gauge}) that Theorem~\ref{thm:barrier} rests on
are fully rigorous and independently verified; the classification claim
above (that any such argument factors through the smoothing-transform
model) can be stated unconditionally as ``standard machinery under
[hypotheses], detailed verification pending'' rather than as
independently re-derived against our exact setup. We never assert that
the residual gauge dispersion $\sigma_Y$ vanishes for the true
arithmetic object; only that the recursion alone cannot force it to.
The missing ingredient that closes the gap \emph{in the idealized
model} --- and which remains to be transferred to the arithmetic
setting --- is cancellation of the centered fine-frequency covariance
after independently averaging the two path indices, together with
equidistribution of the Archimedean phase along paths. Literal
near-independence of paired leaf digits is false by
Theorem~\ref{thm:fresh-digit-coupling}. The known
rigorous analogue, in the forward direction, is the Fourier decay of
Tao's Syracuse variables (Proposition~\ref{prop:tao117}); the
reverse-tree/density version remains open and is \emph{not} implied by
it.
\end{remark}

\subsection{Precise regime of the tail exponent}
\label{sec:endogeny-addendum}

A directed literature search (\S\ref{sec:context}) raised the question
whether $\alpha^\ast=2$ coincides with the \emph{critical case} of the
classical multivariate smoothing-transform theory
\cite{KoleskoMentemeier2015}, characterized by $m'(\alpha)=0$ at a root
of $m(\alpha)=1$ (a boundary/degenerate case in the sense of
Biggins--Kyprianou). Writing $m(\alpha):=q^{\alpha-1}/(2^\alpha-1)$
(the pressure function of \S\ref{sec:pressure} at $q=3$), we compute
\[
  m'(1) = \ln 3 - 2\ln 2\cdot\frac{2^1}{2^1-1}
    \quad\text{explicitly } \approx -0.288 < 0, \qquad
  m'(2) = \ln 3 - \tfrac{4}{3}\ln 2 \approx +0.174 > 0.
\]
Since $m$ is log-convex (as a sum, after clearing denominators, of
exponentials in $\alpha$), two distinct roots of $m(\alpha)=1$
necessarily straddle the unique minimum of $m$, so \emph{neither} can
be a critical (double) root; criticality would require the two roots
to merge, which does not happen at $q=3$. The correct classification
of $\alpha^\ast=2$ is therefore as the \emph{second root with
$m'>0$} --- a Goldie/Guivarc'h-type regularly-varying tail regime
\cite{Goldie1991,BuraczewskiDamekMikosch2016}, not the critical case.
This distinction matters operationally: the critical-case uniqueness
theorem of \cite{KoleskoMentemeier2015} \emph{presupposes} exactly the
independence between subtree copies that is the arithmetic gap of
Theorem~\ref{thm:barrier} --- so even setting aside the regime
mismatch, it could not have supplied a new tool for closing that gap.
We note, tying this to \S\ref{sec:freezing}, that ``second root with
$m'>0$'' is exactly the freezing condition of
Proposition~\ref{prop:always-frozen}: the Goldie/Guivarc'h regime
name correctly classifies the algebraic root $\alpha^\ast=2$ of
$m(\alpha)=1$, but by itself says nothing about whether this root
governs the \emph{quenched} tail of our specific arithmetic recursion
--- which is exactly the question Proposition~\ref{prop:always-frozen}
and Conjecture~\ref{conj:tail-index} address directly.

\section{Empirical calibration of the residual coupling}
\label{sec:calibration}

Proposition~\ref{prop:gauge} shows the gap is real in principle; it
does not quantify how large the residual coupling between sibling
subtrees actually is for the true arithmetic tree. We designed a
three-arm control experiment to calibrate this directly.

\subsection{Design}

\begin{enumerate}[label=Arm~\arabic*.]
\item \textbf{Synthetic i.i.d.\ control}: a Monte Carlo model with the
  exact branch weights $q\cdot 2^{-a}$ but genuinely independent
  subtree contents, truncated by total population size (not depth) to
  mirror the true enumeration procedure.
\item \textbf{Synthetic coupled control}: the same model, but with
  sibling subtree \emph{contents} coupled at adjustable intensity
  $\rho\in[0,1]$ via whole-family replication addressed by a hash of
  (phase, replicate id, tag), constructed so that all single-subtree
  marginals remain exactly unchanged for any $\rho$.
\item \textbf{Real arithmetic tree}: exact enumeration of the true
  reverse tree, re-measured at matched truncation depth.
\end{enumerate}

The residual variance excess, $(\text{Arm 3}) - (\text{Arm 1})$,
compared against the calibration curve $\text{floor}(\rho)$ built from
Arm 2, converts a qualitative statement (``the gap exists'') into a
quantitative bound on the effective coupling $\rho_{\mathrm{eff}}$ that
would be needed in the synthetic model to reproduce whatever residual
is observed in the real tree.

\subsection{Result}

Across the full truncation grid $m=8,\dots,29$, Arm~1 and Arm~3 showed
nearly identical residual variances, with no sign of positive coupling.
A final slope-and-bootstrap measurement at the deepest, best-calibrated
truncation ($m=20$, $4000$ replicate roots, $\rho\in\{0,0.05,0.1,0.2\}$)
gives:
\begin{equation}\label{eq:rho-eff}
  \rho_{\mathrm{eff}} \lesssim 0.06 \quad (95\%\ \text{CI}, \ m=20).
\end{equation}
The excess (real $-$ synthetic) converges monotonically to zero with
increasing truncation depth, the opposite sign from what a genuine
positive arithmetic coupling would produce, reinforcing rather than
merely failing to contradict the conclusion. We state the scope of
this bound precisely: no coupling of the specific functional form
explored by Arm~2 (whole-family replication of subtree \emph{contents}
at intensity $\rho$) was detected in the real tree beyond
$\rho_{\mathrm{eff}}\lesssim0.06$ of that family's maximum intensity.
This is not, by itself, a bound on every conceivable form of coupling;
in particular it does not bound the exact, coarse, $\Delta$-dependent
pair correlation of Proposition~\ref{prop:exact-endogeny} below, which
lives on a disjoint axis (root residue, not subtree content) already
present identically, by construction, in Arm~1 itself --- so it is not
a form of coupling this calibration could have detected or missed. See
Remark~\ref{rem:calibration-consistency} for the verified structural
argument. Equation~\eqref{eq:rho-eff}
converts the purely theoretical gap of Theorem~\ref{thm:barrier} into a
quantitative, falsifiable empirical bound on its size, for this family
of couplings.

\section{A finite-window comparison of the Kontorovich--Lagarias and
Volkov predictions for $5x+1$}
\label{sec:kl}

Independently of the theoretical program above, we report an exact
finite enumeration with a statistical extrapolation.

\subsection{The dispute}

For the reverse tree of $T_5$, Kontorovich--Lagarias
\cite{KontorovichLagarias2009} predict a counting exponent
$\eta_{5,BP}\approx 0.650919$ (identical, as noted in
Remark~\ref{rem:novelty-109}, to our pressure-equation root
$\alpha_-(5)$), while a competing stochastic model due to Volkov,
discussed in the same reference, predicts $\eta^\ast_{5,BP}\approx
0.678$. The authors of \cite{KontorovichLagarias2009} themselves flag
this as an open dispute, unresolved for lack of sufficiently powerful
data; the two predictions differ by only $\Delta=0.027$.

\subsection{Method}

We implemented the exact reverse tree of $T_5$ from scratch (the
admissibility rule of \S\ref{sec:setup} specialized to $q=5$, $d=4$: a
child at exponent $a$ exists from parent $u$ iff $a\equiv
A_0(u\bmod 5)\pmod 4$, with $A_0=\{1\!\mapsto\!4,\allowbreak
2\!\mapsto\!3,\allowbreak 3\!\mapsto\!1,\allowbreak 4\!\mapsto\!2\}$;
nodes with $u\equiv 0\pmod 5$ are sterile). Unlike $q=3$, there is no
parity condition beyond integrality mod $5$.

A first pilot ($60$ roots) already gave a slope of $0.6433$ (95\% CI
$[0.6327,0.6531]$), excluding $0.678$ by roughly $6.4$ standard errors.
A larger production run ($n=300$ roots) initially sampled roots too
close to the measurement window's lower checkpoint, contaminating the
truncation buffer; correcting this (roots restricted well below the
window) gave a slope of $0.63801$ (95\% CI $[0.63159, 0.64410]$) ---
which, at this sample size, excluded \emph{both} candidate values
($0.678$ by $12.6$ standard errors, but also $0.650919$ by $4.1$
standard errors), a failure mode anticipated in advance: once the
standard error falls below the residual truncation bias, a naive
confidence interval on the biased estimator excludes everything,
which is not evidence against the correct prediction.

We therefore rewrote the enumeration to track, in a single pass, the
running path-maximum at every node, giving simultaneous counts at
truncation buffers $10^9$ through $10^{13}$ (validated byte-for-byte
against the earlier buffer-by-buffer method on four test roots). The
pooled mean curve across roots decays geometrically:
\[
  0.60049\ (10^9) \to 0.62387\ (10^{10}) \to 0.63261\ (10^{11}) \to
  0.63650\ (10^{12}) \to 0.63801\ (10^{13}),
\]
with increments $0.0234, 0.0087, 0.0039, 0.0015$ shrinking by a factor
of roughly $0.4$--$0.5$ per decade of truncation --- a signature of a
genuine, extrapolable truncation bias, not noise.

\subsection{Result}

Aitken's $\Delta^2$ extrapolation of this sequence to infinite
truncation gives
\begin{equation}\label{eq:kl-result}
  \hat\eta_5 = 0.639, \qquad 95\%\ \text{bootstrap interval, fixed window} =
  [0.633,\, 0.645].
\end{equation}
The bootstrap interval measures variation between sampled roots after
the truncation extrapolation. It does not include the remaining bias
from the fixed measurement window and therefore cannot exclude either
asymptotic prediction. That bias has a visible signature: the
slope-per-decade panel
\emph{within} the fixed measurement window ($10^4\!\to\!10^5$ through
$10^7\!\to\!10^8$) is still monotonically rising at the last decade
tested ($0.6021\to0.6296\to0.6432\to0.6460$, no plateau), approaching
$0.6509$ monotonically without having reached it --- the signature of
\emph{fixed-window pre-asymptotics} (the window is not yet ``deep
enough''), not of uncorrected truncation bias (already handled by the
Aitken $\Delta^2$ extrapolation above).

\begin{empirical}[Finite-window comparison of two exponent predictions]
\label{thm:kl}
The fixed-window estimator of the $5x+1$ reverse-tree counting exponent,
computed by exact enumeration with Aitken-$\Delta^2$-extrapolated
truncation correction, is
$0.639$ (bootstrap interval $[0.633,0.645]$). The sequence of local
slopes rises monotonically toward the Kontorovich--Lagarias prediction
$0.650919$ on the available windows. Since it has not stabilized, the
experiment does not provide a calibrated confidence interval for the
asymptotic exponent and does not exclude the Volkov prediction.
\end{empirical}

\subsection{A calibrated comparison}

The estimator behind Empirical Result~\ref{thm:kl} had never been run on
a process with a known exponent. It is biased: on a branching random
walk with the same offspring law but an independently drawn branch
class at every node (exponent provably $0.650919$, the pressure-equation
root of \S\ref{sec:pressure}), the same window, truncation, and
extrapolation reads $0.6131$, an under-read of $0.038$, larger than the
gap $\Delta=0.027$ separating the two disputed values. Comparing the raw
reading against either theoretical prediction therefore settles nothing.

The bias depends on how much a process fluctuates, and it can be
calibrated out by comparison rather than subtraction: run the same
estimator on synthetic processes built to have exponent $0.650919$ and
exponent $0.678$ respectively, sharing every feature of the arithmetic
tree except the ultimate value of $q$, and see which reading the
arithmetic tree matches. A sibling congruence in the reverse tree
($w_{j+1}\equiv w_j+(2^d-1)/q\pmod q$ for consecutive children of a
fertile node, a one-line consequence of the recursion of
\S\ref{sec:setup}) makes this construction exact: replacing the integer
denominator $q$ by a tunable real value reproduces the arithmetic
tree's branching and its sibling spacing while moving the exponent to
any target solving $q_{\mathrm{val}}^{\alpha}=q(2^{\alpha}-1)$.

\begin{empirical}[Calibrated separation from the Volkov exponent]
\label{thm:kl-calibrated}
At truncation depth $10^9\!\to\!10^{10}$, where the calibrated controls
read their target exponent to within $0.003$, the arithmetic tree reads
$0.64926$ ($95\%$ interval $[0.64818,0.65027]$); a control built to
exponent $0.650919$ reads $0.64981$ and $0.65122$ on two independent
constructions; a control built to exponent $0.678$ reads $0.67748$,
ten interval-widths from the arithmetic reading. The separation widens,
not narrows, as bias is calibrated away with depth. This measurement
compares against the exponent value $0.678$, not against Volkov's
branching model itself, which has a different tree structure and is not
implemented here.
\end{empirical}

This result stands entirely on its own, independent of the theoretical
program of \S\S\ref{sec:endogeny}--\ref{sec:lemmas}; we highlight it
early in the paper's overall argument because it is the least
framing-dependent contribution we report.

\section{Direct computational test of the Weak Covering Conjecture}
\label{sec:wcc}

\subsection{Wirsching's conjecture}

Wirsching \cite{Wirsching1998} identifies the following combinatorial
covering property as sufficient for uniform positive predecessor
density (his Corollary~II.5.8 and Chapter~V). Define
\[
  R_{j,k} := \Bigl\{ \textstyle\sum_{i=0}^j 2^{\alpha_i} 3^i \;:\;
  j+k\ge \alpha_0>\alpha_1>\cdots>\alpha_j\ge 0 \Bigr\} \subset
  \mathbb{Z}_3^\ast, \qquad |R_{j,k}| = \binom{j+k+1}{j+1}.
\]
Every element of $R_{j,k}$ is automatically a unit mod $3$ (its
top-order term is not divisible by $3$, every other term is), so
``covering $\mathbb{Z}_3^\ast$ mod $3^\ell$'' means hitting all
$2\cdot 3^{\ell-1}$ invertible residue classes.

\begin{conjecture}[Weak Covering Conjecture; Wirsching 1998, Conj.~3.9]
\label{conj:wcc}
There is a sub-exponential function $K(\ell)>0$ (i.e.\
$\lim_\ell K(\ell)e^{-\gamma\ell}=0$ for every $\gamma>0$) such that,
for all $j,\ell$: $|R_{j-1,j}|\ge K(\ell)\cdot 3^\ell \Rightarrow
R_{j-1,j}$ covers $\mathbb{Z}_3^\ast \bmod 3^\ell$.
\end{conjecture}

For $j\ge \ell$, terms of index $i\ge\ell$ vanish mod $3^\ell$, giving
the useful reduction $R_{j-1,j}\bmod 3^\ell = 2^{j-\ell}\cdot
R_{\ell-1,j}\bmod 3^\ell$, which cuts the enumeration cost from
$\binom{2j}{j}$ to $\binom{j+\ell}{\ell}$.

\subsection{Implementation and validation}

We implemented a cyclic-rotation bitset dynamic program: processing
candidate exponents in decreasing order, we maintain, for each count
$c$ of terms chosen so far, the set of reachable partial sums mod
$3^\ell$ as a single Python integer bitset; choosing the exponent $v$
at position $c$ is a cyclic rotation by $2^v\cdot 3^c\bmod 3^\ell$, and
population count of the final bitset gives the image size directly.
The implementation matched an independently computed reference table
exactly ($j^\ast(\ell)$ for $\ell=1,\dots,7$: $1,4,6,7,9,10,11$) and a
pointwise check ($\ell=2,j=2$: image $=\{1,2,5,7\}\bmod 9$, missing
$\{4,8\}$).

\subsection{Results}

We computed $j^\ast(\ell)$, the smallest $j$ for which $R_{j-1,j}$
covers, for $\ell=1,\dots,20$ (cost grew by a factor $\approx 3.3$ per
step near the end; $\ell=20$ alone took $\approx 27$ minutes; further
extension in pure Python was judged not worthwhile). Writing
$e(\ell):=j^\ast(\ell)-\log_4 3\cdot\ell$ for the excess over the
entropy-matched threshold $\log_4 3\approx 0.7925$:

\begin{center}
\begin{tabular}{c|c|c|c}
$\ell$ & $j^\ast$ & $j^\ast/\ell$ & $e(\ell)$ \\ \hline
1 & 1 & 1.0000 & 0.208 \\
5 & 9 & 1.8000 & 5.038 \\
10 & 15 & 1.5000 & 7.075 \\
15 & 20 & 1.3333 & 8.113 \\
16 & 20 & 1.2500 & 7.320 \\
17 & 21 & 1.2353 & 7.528 \\
18 & 22 & 1.2222 & 7.735 \\
19 & 23 & 1.2105 & 7.943 \\
20 & 24 & 1.2000 & 8.150
\end{tabular}
\end{center}

(Full table for all $20$ values is in the accompanying computational
record.) $j^\ast(\ell)$ exists and is finite throughout the tested
range. An earlier reading of the first $17$ points suggested a
monotonically decelerating local slope approaching the threshold
$\log_4 3$, which the three new points at $\ell=18,19,20$ (all
increments exactly $+1$, no new plateaus) revealed to be an artifact
of an isolated plateau at $\ell=16$ crossing the sliding measurement
window: without it, the local slope oscillates in $0.74$--$0.86$,
straddling $\log_4 3$ from both sides rather than converging cleanly.

Model comparison on the tail $\ell=10$--$20$ (constant vs.\ logarithmic
vs.\ square-root vs.\ slow-linear growth of $e(\ell)$, by AIC) now
disfavors pure stabilization ($\Delta\mathrm{AIC}\approx 5$), consistent
with a plateau-frequency test (only $1$ plateau observed in $19$
increments versus $\approx 3.9$ expected under stabilization at rate
$\log_4 3$; binomial $p\approx 0.072$, marginal but in the same
direction). The four slow-growth models remain statistically
indistinguishable from one another in this range
($\Delta\mathrm{AIC}<0.5$); they would only separate by $\ge 1$ unit of
$j^\ast$ near $\ell\approx 40$, out of reach of this algorithmic
approach.

\begin{empirical}[Direct computational test of the Weak Covering
Conjecture]\label{thm:wcc}
$j^\ast(\ell)$ exists and is finite for all $\ell\le 20$ (the first
direct computation here of the exact object of
Conjecture~\ref{conj:wcc}). The excess $e(\ell)$ qualitatively favors
unbounded slow growth over stabilization (two independent statistics,
$\Delta\mathrm{AIC}$ and plateau frequency, in the same direction but
not individually decisive), consistent with the weak form
(Conjecture~\ref{conj:wcc}) and disfavoring the corresponding strong
form with $K(\ell)$ bounded.
\end{empirical}

\section{Three precision regimes for Tao's bivariate extension}
\label{sec:regimes}

\subsection{Why the naive bivariate argument fails}

A natural strategy for closing the gap of Theorem~\ref{thm:barrier} is
to apply Proposition~\ref{prop:tao117} twice, once per sibling leaf,
and combine via some form of conditional independence given the common
ancestor. A close reading of the mechanism of Tao's Section~7 proof
(characters $\chi_\xi$, i.i.d.\ Pascal-distributed valuation pairs, a
phase evolving multiplicatively by $\log 9/\log 2$, a ``black set'' of
well-separated triangles, a two-dimensional renewal process with
monotonicity) shows that this mechanism operates entirely on
$\mathrm{Syrac}(\mathbb{Z}/3^n\mathbb{Z})$ as constructed from
\emph{independent} geometric digits by definition
\eqref{eq:syracuse-def}; the link to actual integers is made elsewhere
in \cite{Tao2022} (via a total-variation approximation valid only for
depth $n\lesssim\log N$).

For our problem --- a single fixed integer $v$, two sibling leaves
$w_1(v),w_2(v)$ --- the two leaves are \emph{not} two i.i.d.\ samples;
they are two different deterministic functions of the \emph{same}
variable $v$.

\begin{proposition}[Perfect dependence at fixed depth]\label{prop:fixed-pair}
For a fixed pair of sibling paths of depth $D$, the existence
congruence pins $v$ to a single class mod $3^D$: writing
$v=v_0+3^D t$, the two leaves become affine functions of the same free
variable $t$, with unit multipliers: $w_i=w_i(v_0)+2^{A_i}t$.
Consequently the pair $(w_1,w_2)\bmod 3^\ell$ lies on a line of size
$3^\ell$ inside $(\mathbb{Z}/3^\ell\mathbb{Z})^2$, for every $\ell\ge 1$
--- perfect deterministic dependence at every precision, with an entire
line of resonant frequencies where the correlation has modulus exactly
$1$.
\end{proposition}

\begin{theorem}[Maximal coupling of fresh digits]
\label{thm:fresh-digit-coupling}
Fix a pair of paths and condition their leaf values modulo $3^r$.
If the free arithmetic parameter is uniform over the $3^s$ compatible
lifts modulo $3^{r+s}$, let $X_s,Y_s$ denote the next $s$ base-$3$
digits of the two leaves. Then both marginals are uniform and
\[
 d_{\rm TV}\bigl(\mathcal L(X_s,Y_s),
     \mathcal L(X_s)\otimes\mathcal L(Y_s)\bigr)=1-3^{-s},
 \qquad
 I(X_s;Y_s)=s\log3.
\]
Thus fresh digit blocks of a fixed leaf pair become maximally, rather
than weakly, dependent as $s\to\infty$.
\end{theorem}

\begin{proof}
By Proposition~\ref{prop:fixed-pair}, after the coarse class is fixed
the two leaves are affine functions of one lift parameter, and both
multipliers are units modulo $3^s$. The map from the fresh block of the
first leaf to that of the second is therefore a permutation of
$3^s$ points. The joint law is uniform on its graph, while the product
law is uniform on all $3^{2s}$ pairs. Direct summation gives total
variation $1-3^{-s}$. Since one block determines the other, their
mutual information equals the entropy of a uniform $3^s$-point
variable, namely $s\log3$.
\end{proof}

Both identities were checked numerically in
\texttt{collatz-endogeny/\allowbreak sec8-fresh-digit-coupling}.

Consequently, ``apply Proposition~\ref{prop:tao117} twice, assuming
conditional independence given the ancestor'' is not merely circular:
it is false as stated, at any precision.

\subsection{The correct, non-circular reformulation}

The quantity that actually matters --- decorrelation of \emph{subtree
aggregates}, not leaf-to-leaf --- is a second-moment statement,
expanding into a sum over \emph{pairs of paths}:
\[
  \mathbb{E}_v[Z_1 Z_2] \;\propto\; \#\{(a,a') : a_1\in B_1,\ a'_1\in
  B_2,\ \mathrm{Syrac}(a)\equiv\mathrm{Syrac}(a') \bmod 3^D\},
\]
which by Fourier inversion gives $\Cov \propto \sum_{\xi\ne 0}
S_1(\xi)S_2(\xi)^\ast$, with $S_i$ character sums conditioned on the
first step. Here there is no circularity: the two path indices are
independent by the tautology of expanding a square, not by an
arithmetic hypothesis.

\subsection{The three regimes}

\begin{enumerate}
\item \textbf{Fixed pair, any $\ell$}: perfect dependence
  (Proposition~\ref{prop:fixed-pair}); no decay is possible, and
  decorrelation can only come from averaging over paths, never
  leaf-to-leaf.
\item \textbf{$\ell = O(\log D)$}: \emph{provable}, with Tao's
  Section~7 machinery essentially unchanged. The budget
  $|\rho(\xi)|\le C_A D^{-A}$ closes the sum over $3^\ell-1$
  frequencies provided $3^\ell \ll D^{2A}$, i.e.\ polynomial moduli in
  $D$ --- an honest, reachable lemma, though its most natural
  formulation turns out to be false as originally sketched
  (\S\ref{sec:lemmas}).
\item \textbf{$\ell \asymp c\cdot D$} (where the endogeny barrier, the
  fresh digits, and Wirsching's Weak Covering Conjecture all live):
  this requires exponential control of the linear-precision
  frequency tail after sublinear-conductor modes have been separated;
  Proposition~\ref{prop:tao117} gives only
  super-polynomial decay, and Tao's Section~7 method structurally
  cannot give more, because the losses in the black-set argument are
  tied to Diophantine properties of $\log 3/\log 2$ via the slope
  $\log 9/\log 2$.
\end{enumerate}

\begin{remark}[Nearby open theories]\label{thm:regime3}
Regime 3 has features in common with correlations of digital functions
in bases $2$ and $3$ \cite{Spiegelhofer2021}, effective rigidity for
$\times2,\times3$ actions, and Fourier decay for self-similar measures.
We do not prove a reduction to or from any of these problems. The target
used below is an exponential bound for the linear-precision part of the
conditioned character sums $S_i(\xi)$, with sublinear-conductor modes
treated through Theorem~\ref{thm:sublinear-precision-ensemble}.
\end{remark}

\section{Triangulation: connections among three formulations of the
obstruction}
\label{sec:triangulation}

\subsection{$\beta=1$ (Tao 2020) and a weighted covering problem}

Tao \cite{Tao2020} defines $c_n := \inf_{b\in\mathbb{Z}/3^n\mathbb{Z},\,
3\nmid b} \Pr(\mathrm{Syrac}(\mathbb{Z}/3^n\mathbb{Z})=b)$ and proposes
the conjecture $\beta=1$: $c_n = 3^{-n+o(n)}$ (``no residue class has
probability much smaller than average''), proving conditionally that
$\beta=1$ implies Collatz density $\gg x^{1-o(1)}$ (almost total), and
noting that the best known unconditional bound is Krasikov--Lagarias'
$\gamma=0.84$ \cite{KrasikovLagarias2003}.

\begin{proposition}[Support identity and geometric cost]\label{prop:beta-wcc}
After a unit twist, Tao's Syracuse random variable and Wirsching's sets
$R_{j,k}$ are the same algebraic object (sums of powers of $2$ and $3$
with strictly decreasing exponents). Let $B_\ell(z)$ be the least total
geometric cost $a_1+\cdots+a_\ell$ of a Syracuse tuple representing
$z\pmod {3^\ell}$, and let $j^\ast(\ell)$ be the least $j$ for which
$R_{\ell-1,j}$ covers every admissible residue. Then
\[
  \max_z B_\ell(z)=\ell+j^\ast(\ell).
\]
More specifically, tuples of cost at most $\ell+j$ are in bijection,
after multiplication by $2^{\ell+j}$, with $R_{\ell-1,j}$ modulo
$3^\ell$.
\end{proposition}

Ordinary covering controls the cheapest representation of each residue,
whereas $c_\ell$ sums the geometric weights of all its representations.
At the WCC threshold $j^\ast(\ell)=\log_4(3)\ell+o(\ell)$, the cheapest
representation alone gives
\[
 c_\ell\ge 2^{-\ell-j^\ast(\ell)}
 =3^{-(1/2+\log_3 2)\ell+o(\ell)}.
\]
This yields only $\beta\le 1/2+\log_3 2=1.1309\ldots$, rather than
$\beta=1$ (identity and bound checked against brute force in
\texttt{collatz-endogeny/\allowbreak sec9-wcc-beta-bridge}). The
entropy implication previously asserted here is invalid.
The missing statement is a weighted covering estimate reaching the
typical cost window: the total geometric multiplicity must be within a
subexponential factor of its mean uniformly over residues. Our
computation of $j^\ast(\ell)$ tests WCC, but does not by itself test
this weighted assertion.

\begin{theorem}[Large-deviation barrier at the WCC cost]
\label{thm:wcc-large-deviation}
Let $A_1,\ldots,A_\ell$ be independent with
$\Pr(A_i=m)=2^{-m}$ for $m\ge1$, and put $B_\ell=\sum_iA_i$. For
$1<s<2$,
\begin{equation}\label{eq:wcc-rate}
 \Pr(B_\ell\le s\ell)
 =\exp\{-I(s)\ell+o(\ell)\},
 \qquad
 I(s)=s\log2+(s-1)\log(s-1)-s\log s>0.
\end{equation}
Let $\mu_\ell^{(\le s)}$ be the subprobability Syracuse measure formed
from tuples with $B_\ell\le s\ell$. Then, with the minimum over unit
residues,
\[
 \min_a\mu_\ell^{(\le s)}(a)
 \le 3^{-\ell}\exp\{-I(s)\ell+o(\ell)\}.
\]
At $s_\ast=1+\log_4 3$, the rate is
$I(s_\ast)=0.0120393866\ldots$, and the corresponding exponent relative
to base $3$ is
\[
 1+\frac{I(s_\ast)}{\log3}=1.0109587219\ldots>1.
\]
Thus even perfect weighted equidistribution inside the critical WCC
cost slice cannot establish $\beta=1$.
\end{theorem}

\begin{proof}
The exact negative-binomial formula is
\[
 \Pr(B_\ell=m)={m-1\choose\ell-1}2^{-m},\qquad m\ge\ell.
\]
For $m<2\ell-2$ these masses increase with $m$. Hence the logarithm of
their sum through $\lfloor s\ell\rfloor$ differs from the logarithm of
the final term by at most $O(\log\ell)$. Stirling's formula applied to
that term gives \eqref{eq:wcc-rate}. The total mass of
$\mu_\ell^{(\le s)}$ is the probability in \eqref{eq:wcc-rate}. Its
minimum over the $2\cdot3^{\ell-1}$ unit residues is at most its
average. The fixed positive factor between this average and
$3^{-\ell}$ is absorbed by $o(\ell)$. Direct substitution of
$s_\ast$ gives the stated constants.
\end{proof}

The finite-level rate and bound are computed in
\texttt{collatz-endogeny/\allowbreak sec9-wcc-beta-bridge}.

Theorem~\ref{thm:wcc-large-deviation} sharpens the distinction between
O2 and O3. Multiplicity within the WCC budget can improve the
single-representation exponent $1.1309\ldots$ only as far as the
average truncated exponent $1.01096\ldots$. A route to $\beta=1$ must
reach costs $B_\ell=2\ell-o(\ell)$, where the cost law has
subexponential mass. This is the central microcanonical regime of
Theorem~\ref{thm:microcanonical}.

Tao's $2020$ approach also shows that the endogeny/decorrelation route
of \S\ref{sec:endogeny} is not the only path: it replaces
``near-independence between subtrees'' with ``deterministic separation
of preimages (easy) plus $L^\infty$ marginal control ($\beta=1$, where
all the difficulty lives)''. This confirms that a route avoiding
decorrelation altogether exists, and that it hits the \emph{same} wall
in marginal form.

\subsection{Wirsching's 2003 functional equation and the base-$3$
Fabius function}

A different, more refined route toward the same goal is
Wirsching~\cite{Wirsching2003}, whose central combinatorial object,
the ``Elka'' functions $e_\ell(k,a):=|E_{\ell,k}(a)|$, count paths of
length $k+\ell$ in the Collatz graph ending at $a$ --- the same
counting object underlying our entire $G(v)/\mathrm{Syrac}$ program,
now viewed on the state space $\mathbb{X}=\mathbb{I}_3\times
\mathbb{Z}_3^\times$: the $\mathbb{Z}_3^\times$ factor carries the
$3$-adic variable (habitat of the Syracuse measure), while the
$\mathbb{I}_3\cong\{0,1,2\}^{\mathbb{N}}$ factor carries the normalized
step-count $k/3^\ell$, converging to a density $\varphi$.

\begin{proposition}[Base-$3$ Fabius function]\label{prop:fabius}
$\varphi$ is the density of $X = \sum_{j\ge 1} 2U_j\cdot 3^{-j}$, $U_j$
i.i.d.\ uniform on $[0,1]$ --- the base-$3$ analogue of the Fabius
function (which satisfies $f'(x)=2f(2x)$; here $\varphi'(x) =
\tfrac{9}{2}\varphi(3x)$ on $[0,2/3]$), a member of the family of
``atomic functions'' of Rvachev. It is the unique $L^1([0,1])$ fixed
point of the averaging operator $W_3 f(x) := \tfrac{3}{2}
\int_{3x-2}^{3x} f(t)\,dt$, $C^\infty$, piecewise polynomial away from
the standard Cantor set, with certified geometric convergence
$\|f_n-\varphi\|_1\le 2^{-n+1}\|f_1-f_0\|_1$.
\end{proposition}

Wirsching's chain of implications reduces uniform positive predecessor
density to three proposed steps. The first has an elementary proof.

\begin{theorem}[Wirsching's Conjecture 1]\label{thm:wirsching-conj1}
In the notation of \cite{Wirsching2003}, condition $(?2)$ implies
condition $(?1)$, with a possibly smaller central-limit window.
\end{theorem}

\begin{proof}
Put $c_0=1$ and $c_j=2\cdot3^{j-1}$ for $j\ge1$. If $q_\ell(k)$ counts
the bounded urn distributions defining $\bar g_\ell(k)$, then
\[
 \sum_{k\ge0}q_\ell(k)z^k
 =\prod_{j=0}^\ell\frac{1-z^{c_j}}{1-z},\qquad
 \sum_{m\ge0}p_\ell(m)z^m
 =\prod_{j=0}^\ell(1-z^{c_j})^{-1}.
\]
Since $q_\ell(k)=2\cdot3^{\ell-1}\bar g_\ell(k)$ and
$\bar e_\ell=p_\ell*\bar g_\ell$, cancellation gives
\[
 \bar e_\ell(k)=\frac{1}{2\cdot3^{\ell-1}}
                 {k+\ell\choose\ell}.
\]
The unrestricted partition count for the coins $1,2,6,18,\ldots$ is
at most $\exp(C\log^2(m+2))$. Also
\[
 \frac{{k-m+\ell\choose\ell}}{{k+\ell\choose\ell}}
 \le \left(\frac{k}{k+\ell}\right)^m.
\]
Uniformly for $|k-\ell|\le\delta\sqrt\ell$, the part of the convolution
with $m\ge\eta\sqrt\ell$ is
$\exp(-c\sqrt\ell+O(\log^2\ell))$. Choose
$\delta+\eta<\delta_1$. On the remaining terms $j=k-m$ lies in the
window of $(?2)$, and positivity transfers its lower bound from
$g_\ell/\bar g_\ell$ to $e_\ell/\bar e_\ell$.
\end{proof}

\begin{remark}[Relation to the source]\label{rem:wirsching-conj1-source}
Wirsching records the two convolution identities used above and states
the implication as Conjecture~1 \cite[Sections~1--2]{Wirsching2003}.
The annotated Collatz bibliography also lists the three implications
in that paper as conjectures \cite{LagariasBibliographyII}. The proof
above supplies the generating-function cancellation and the
convolution-tail estimate for the first implication. The generating-function
identity is checked coefficient by coefficient in
\texttt{collatz-endogeny/\allowbreak sec9-wirsching-conjecture1}.
\end{remark}

The second conjecture remains open. The most concrete sufficient
condition above it is Conjecture~3, an asymptotic assertion on
$\varphi(z_\ell)/\varphi_0(z_\ell)$ along the central-limit window
$|\ell-k_\ell|\le\delta\sqrt\ell$, which we tested numerically.

\subsection{Certified numerical test of Wirsching's Conjecture 3}

Using the exact self-similarity $X\overset{d}{=}(2U+X)/3$, the moments
$M_i=\mathbb{E}[X^i]$ satisfy the exact rational recursion
\[
  M_i = \frac{1}{3^i-1}\sum_{k=1}^i \binom{i}{k}\frac{2^k}{k+1} M_{i-k},
\]
and, combined with an antiderivative reduction expressing $\varphi$ at
extreme-tail arguments in terms of exact binomial moment sums (rather
than iterating $W_3$, whose $L^1$ error certificate is too weak to
control pointwise values at the relevant scale), we obtain a
zero-truncation-error evaluation of $\varphi$ at depth $\ell$ up to
$\ell=500$ (working throughout in log-space with $60$-digit precision;
error certified $\le 10^{-8}$).

\begin{empirical}[Numerical test of Wirsching's Conjecture 3]
\label{thm:conjecture3}
The decisive ratio $L_\ell := 3^{1-\ell}\varphi(3x_\ell^+)/\varphi(x_\ell^+)$,
measured for $\ell=250,\dots,500$ in the CLT window ($u\in[-2,2]$),
converges to the predicted value $2/3$ with deficit
$\ell\cdot(2/3-L_\ell) \to 0.580\pm 0.001$ --- a coefficient
independently reproduced by Berg--Krüppel's own asymptotic
$\varphi_0$. The quantity $\ln(\varphi/\varphi_0)$ is increasing,
concave, and uniform across $u\in[-2,2]$ (dispersion $<10^{-4}$),
consistent with a finite limit $L=-0.619\pm0.001\,(\mathrm{stat})$,
with a tail-shape systematic uncertainty of $\pm 0.015$, giving
$c=e^L\approx 0.538$ (honest range $0.53$--$0.55$), best fit by a
tail term $0.79/\sqrt\ell$; no log-periodic modulation is detected.
This supports Conjecture~3 with certified numerical error on the tested
range. Conjecture~2 remains open; Theorem~\ref{thm:wirsching-conj1}
settles Conjecture~1 independently of this computation.
\end{empirical}

\begin{theorem}[Microcanonical decomposition]\label{thm:microcanonical}
Let $g_\ell(k,a)$ be Wirsching's generator, let
$c_i=2\cdot3^i$, and let $\mu_\ell$ be the Syracuse law modulo
$3^\ell$. For every unit residue $a$,
\begin{equation}\label{eq:microcanonical}
 \mu_\ell(a)=
 \left(\prod_{i=0}^{\ell-1}\frac{1}{2(1-2^{-c_i})}\right)
 \sum_{k=0}^{3^\ell-\ell-1}2^{-k}g_\ell(k,a).
\end{equation}
Suppose there are $\delta,\eta>0$ such that, for all sufficiently large
$\ell$, every unit $a$, and every integer
$|k-\ell|\le\delta\sqrt\ell$,
\[
 g_\ell(k,a)\ge\eta\,\overline g_\ell(k),
 \qquad
 \overline g_\ell(k)=\frac{1}{2\cdot3^{\ell-1}}
 \sum_{b\in(\mathbb Z/3^\ell\mathbb Z)^\times}g_\ell(k,b).
\]
Then $\mu_\ell(a)\ge C_{\delta,\eta}3^{-\ell}$ uniformly in $a$ and
$\ell$. In particular, condition $(?3)$ of \cite{Wirsching2003}
implies the cylinder lower bound required by $\beta=1$.
\end{theorem}

\begin{proof}
At recursion level $i$, multiplication by $2$ has order
$c_i=2\cdot3^i$ modulo $3^{i+1}$. If $A_i$ has the geometric law
$\Pr(A_i=m)=2^{-m}$, then $J_i=A_i-1\bmod c_i$ satisfies
\[
 \Pr(J_i=j)=\frac{2^{-(j+1)}}{1-2^{-c_i}},
 \qquad 0\le j<c_i.
\]
Folding each exponent independently by the applicable order leaves the
Syracuse residue unchanged. Wirsching's recursion for $g_\ell$ counts
the folded vectors with $\sum_iJ_i=k$. Multiplying their probabilities
gives \eqref{eq:microcanonical}.

Put $K_\ell=\sum_{i<\ell}J_i$. The summands are independent and are
uniformly dominated by an untruncated geometric variable. Direct
summation gives
$\mathbb E K_\ell=\ell+O(1)$ and
$\operatorname{Var}K_\ell=2\ell+O(1)$. The Lindeberg central limit
theorem therefore gives
\[
 \Pr(|K_\ell-\ell|\le\delta\sqrt\ell)
 \longrightarrow 2\Phi(\delta/\sqrt2)-1>0.
\]
Sum \eqref{eq:microcanonical} over the central values of $k$ and use
the assumed lower bound. Since there are $2\cdot3^{\ell-1}$ unit
residues, this yields
\[
 \mu_\ell(a)\ge
 \frac{\eta}{2\cdot3^{\ell-1}}
 \Pr(|K_\ell-\ell|\le\delta\sqrt\ell)
 \ge C_{\delta,\eta}3^{-\ell}.
\]
\end{proof}

\begin{proposition}[Conditional local-limit bridge]
\label{prop:complex-deconditioning}
Set
\[
 G_{\ell,a}(z)=\sum_k g_\ell(k,a)z^k,\qquad
 \overline G_\ell(z)=\frac{1}{2\cdot3^{\ell-1}}\sum_aG_{\ell,a}(z).
\]
Define the conditional cost laws
\[
 Q_{\ell,a}(k)=\frac{2^{-k}g_\ell(k,a)}{G_{\ell,a}(1/2)},
 \qquad
 \overline Q_\ell(k)=
 \frac{2^{-k}\overline g_\ell(k)}{\overline G_\ell(1/2)}.
\]
Then the exact identity
\begin{equation}\label{eq:deconditioning}
 \frac{g_\ell(k,a)}{\overline g_\ell(k)}
 =
 \frac{G_{\ell,a}(1/2)}{\overline G_\ell(1/2)}
 \frac{Q_{\ell,a}(k)}{\overline Q_\ell(k)}
\end{equation}
holds. Fix $\delta>0$. If, uniformly in unit $a$ and
$|k-\ell|\le\delta\sqrt\ell$,
\[
 \frac{G_{\ell,a}(1/2)}{\overline G_\ell(1/2)}\ge\eta>0,
 \qquad
 Q_{\ell,a}(k)\ge c_\delta\ell^{-1/2},
\]
then condition $(?3)$ holds.
\end{proposition}

\begin{proof}
Identity \eqref{eq:deconditioning} follows by cancellation. Under the
folded geometric law in the proof of Theorem~\ref{thm:microcanonical},
$\overline Q_\ell$ is the law of $K_\ell$.
The summands have uniformly bounded exponential moments,
mean sum $\ell+O(1)$, and variance $2\ell+O(1)$. Their lattice local
central limit theorem gives, uniformly in the stated window,
\[
 \overline Q_\ell(k)\asymp_\delta\ell^{-1/2}.
\]
The two assumed lower bounds and \eqref{eq:deconditioning} now give a
positive uniform lower bound for
$g_\ell(k,a)/\overline g_\ell(k)$.
\end{proof}

The first factor in \eqref{eq:deconditioning} is the canonical Syracuse
mass relative to its mean. The conjecture $\beta=1$ gives only a
subexponential lower bound for this factor; the proposition assumes a
constant lower bound. The second factor is a local limit theorem after
conditioning on the residue. Its Fourier representation is
\[
 Q_{\ell,a}(k)=\frac{1}{2\pi}\int_{-\pi}^{\pi}
 e^{-ikt}\frac{G_{\ell,a}(e^{it}/2)}{G_{\ell,a}(1/2)}\,dt.
\]
Thus O3 requires phase control of the conditional cost law in addition
to the real-axis information in O2.

\begin{remark}[Local limit theorems for operator cocycles]
Spectral local limit theorems exist for random expanding dynamics
\cite{DragicevicEtAl2018} and for time-dependent expanding systems
\cite{Hafouta2020}. Their hypotheses include a twisted operator
cocycle acting on a fixed Banach space, a uniform spectral or
Ruelle--Perron--Frobenius gap near zero, and an aperiodicity estimate
away from zero. These hypotheses would give the phase control in
Proposition~\ref{prop:complex-deconditioning}. They have not been
verified for Wirsching's $S_\ell$. Theorem~3 of
\cite{Wirsching2003} gives strong convergence on fixed
equicontinuous families, while the densities used here acquire
resolution $3^{-\ell}$. Thus the available local limit theorems do not
remove the growing-resolution obstruction of H-134.
\end{remark}

\begin{theorem}[Equivalence of ensembles at fixed precision]
\label{thm:fixed-precision-ensemble}
For a unit residue $a$ put
\[
 p_{\ell,k}(a)=\frac{g_\ell(k,a)}{\sum_b g_\ell(k,b)}.
\]
Fix $r\ge1$ and $\delta>0$. Uniformly over integers
$|k-\ell|\le\delta\sqrt\ell$, the projection of $p_{\ell,k}$ modulo
$3^r$ converges in total variation to the Syracuse law $\mu_r$ as
$\ell\to\infty$.
\end{theorem}

\begin{proof}
Under the canonical weights in Theorem~\ref{thm:microcanonical}, the
folded costs $J_0,\ldots,J_{\ell-1}$ are independent and
$p_{\ell,k}$ is the law of the resulting residue conditional on
$K_\ell=\sum_iJ_i=k$. Fix the terminal block of $r$ costs. For each
fixed value $\mathbf j$ of that block, with coordinate sum
$s(\mathbf j)$, independence gives
\[
 \Pr(\mathbf J=\mathbf j\mid K_\ell=k)
 =\Pr(\mathbf J=\mathbf j)
   \frac{\Pr(K_{\ell-r}=k-s(\mathbf j))}
        {\Pr(K_\ell=k)}.
\]
The lattice local central limit theorem used in
Proposition~\ref{prop:complex-deconditioning} makes the ratio tend to
one, uniformly in the stated window, for every fixed $\mathbf j$.
The geometric domination of the coordinates makes the block tails
uniformly summable. Splitting at block sum $o(\sqrt\ell)$ and applying
the same local estimate there extends the convergence to total
variation on the whole block.

The terminal $r$ folded costs determine the final residue modulo
$3^r$. At a terminal stage whose original cap is $C$, reduction to the
cap $d=2\cdot3^i$ required at precision $r$ is exact because $d\mid C$
and
\[
 \sum_{\substack{0\le j<C\\j\equiv s\ ({\rm mod}\ d)}}
 \frac{2^{-(j+1)}}{1-2^{-C}}
 =\frac{2^{-(s+1)}}{1-2^{-d}},\qquad 0\le s<d.
\]
The reduced block therefore has exactly the canonical folded-geometric
law defining $\mu_r$. Total variation cannot increase under the
deterministic map from the block to its residue, which proves the
claim.
\end{proof}

\begin{theorem}[Equivalence of ensembles at sublinear precision]
\label{thm:sublinear-precision-ensemble}
Let $r=r_\ell$ be an integer sequence with $1\le r\le\ell$ and
$r=o(\ell)$. Fix $\delta>0$. Uniformly over integers
$|k-\ell|\le\delta\sqrt\ell$, the projection of $p_{\ell,k}$ modulo
$3^r$ has total-variation distance $o(1)$ from $\mu_r$.
\end{theorem}

\begin{proof}
It remains to treat $r\to\infty$, since bounded subsequences follow
from Theorem~\ref{thm:fixed-precision-ensemble}. Split the canonical
cost as $K_\ell=R_{\ell,r}+B_{\ell,r}$, where $B_{\ell,r}$ is the sum
of the terminal $r$ folded costs. The Radon derivative of the law of
the terminal block conditional on $K_\ell=k$, relative to its
unconditioned law, is
\begin{equation}\label{eq:block-radon}
 D_{\ell,r}(B_{\ell,r})
 =\frac{\Pr(R_{\ell,r}=k-B_{\ell,r})}
        {\Pr(K_\ell=k)}.
\end{equation}
Choose $h_\ell\to\infty$ slowly enough that
$h_\ell\sqrt{r/\ell}\to0$. Uniform exponential moments give
\[
 \Pr\bigl(|B_{\ell,r}-\mathbb EB_{\ell,r}|>
                 h_\ell\sqrt r\bigr)=o(1).
\]
On the complementary event, removing the block changes the centered
local-CLT coordinate by $o(1)$ and changes the variance by a relative
$o(1)$. The uniform lattice local central limit theorem for the folded
geometric triangular array therefore gives
$D_{\ell,r}(B_{\ell,r})=1+o(1)$ uniformly there. Since
$\Pr(K_\ell=k)\gg_\delta\ell^{-1/2}$ in the central window, while the
standard concentration bound for lattice sums with uniformly bounded
exponential moments gives
$\sup_m\Pr(R_{\ell,r}=m)\ll(\ell-r)^{-1/2}$. Hence
$D_{\ell,r}\ll_\delta1$ on its whole support. Truncation of the
exceptional set now gives
$\mathbb E|D_{\ell,r}-1|=o(1)$. This is twice the total-variation
distance between the conditioned and unconditioned terminal blocks.
Their deterministic residue maps cannot increase that distance. The
cap-reduction identity in the proof of
Theorem~\ref{thm:fixed-precision-ensemble} shows that the unconditioned
residue has law exactly $\mu_r$.
\end{proof}

\begin{theorem}[Linear-block nonequivalence]\label{thm:linear-block-nonequivalence}
Let $\mathbf B_{\ell,r}=(J_{\ell-r},\ldots,J_{\ell-1})$ be the
terminal block of folded costs and let
$S_{\ell,r}=\sum_{i=\ell-r}^{\ell-1}J_i$. Suppose that
\[
 \frac r\ell\longrightarrow\rho\in(0,1),\qquad
 \frac{k_\ell-\ell}{\sqrt\ell}\longrightarrow u.
\]
Then
\begin{align*}
 &d_{\rm TV}\bigl(
   \mathcal L(\mathbf B_{\ell,r}\mid K_\ell=k_\ell),
   \mathcal L(\mathbf B_{\ell,r})\bigr)\\
 &\quad=d_{\rm TV}\bigl(
   \mathcal L(S_{\ell,r}\mid K_\ell=k_\ell),
   \mathcal L(S_{\ell,r})\bigr)\\
 &\quad\longrightarrow
 d_{\rm TV}\bigl(
   \mathcal N(\rho u,2\rho(1-\rho)),
   \mathcal N(0,2\rho)\bigr)>0,
\end{align*}
where the second parameter of $\mathcal N$ is the variance. In fact,
the first equality holds without the $o(1)$ at every finite level.
For $u=0$, the limit equals
\[
 2\left[
 \Phi\left(\sqrt{\frac{-\log(1-\rho)}{\rho}}\right)
 -\Phi\left(\sqrt{\frac{-(1-\rho)\log(1-\rho)}{\rho}}\right)
 \right].
\]
\end{theorem}

\begin{proof}
Write $R_{\ell,r}=K_\ell-S_{\ell,r}$. Relative to the
unconditioned block law, the conditioned law has density
\[
 \frac{\Pr(R_{\ell,r}=k_\ell-S_{\ell,r})}
      {\Pr(K_\ell=k_\ell)}.
\]
This density is a function of $S_{\ell,r}$ alone. Taking its absolute
deviation from one proves the finite-level equality of the two total
variation distances.

The folded costs have uniformly bounded exponential moments. Their
means and variances differ from $1$ and $2$, respectively, by summable
errors. The lattice local central limit theorem, applied separately to
$S_{\ell,r}$ and $R_{\ell,r}$, gives the unconditioned limit
$\mathcal N(0,2\rho)$ for
$(S_{\ell,r}-r)/\sqrt\ell$. After conditioning, its limiting density is
\[
 x\longmapsto
 \frac{\phi_{0,2\rho}(x)\phi_{0,2(1-\rho)}(u-x)}
      {\phi_{0,2}(u)}
 =\phi_{\rho u,2\rho(1-\rho)}(x).
\]
Uniform exponential tails promote the local convergence to convergence
in total variation. The two limiting variances are different, so the
distance is positive. When $u=0$, the two centered Gaussian densities
cross at
$|x|=\sqrt{-2(1-\rho)\log(1-\rho)}$; integration on the central
interval gives the displayed formula.
\end{proof}

Theorem~\ref{thm:linear-block-nonequivalence} shows that the condition
$r=o(\ell)$ in Theorem~\ref{thm:sublinear-precision-ensemble} is sharp
for the latent cost vectors. It does not prove nonequivalence after
projection modulo $3^r$: the residue map can discard the block sum,
and total variation can decrease under that projection. The
linear-precision residue and Fourier problems remain open.

\begin{theorem}[Full-precision likelihood domination]
\label{thm:ensemble-divergence}
Let $p_{\ell,k}^{(r)}$ be the projection of $p_{\ell,k}$ modulo $3^r$,
where $1\le r\le\ell$. For every set
$E\subseteq\mathbb Z/3^r\mathbb Z$,
\begin{equation}\label{eq:ensemble-domination}
 p_{\ell,k}^{(r)}(E)
 \le \Pr(K_\ell=k)^{-1}\mu_r(E).
\end{equation}
Consequently
\[
 D_{\rm KL}(p_{\ell,k}^{(r)}\Vert\mu_r)
 \le D_\infty(p_{\ell,k}^{(r)}\Vert\mu_r)
 \le-\log\Pr(K_\ell=k).
\]
Uniformly for $|k-\ell|\le\delta\sqrt\ell$, the last quantity is
$\frac12\log\ell+O_\delta(1)$.
\end{theorem}

\begin{proof}
On the canonical folded-cost space, the microcanonical law is the
conditional law given $K_\ell=k$. Hence, for the inverse image
$\widetilde E$ of $E$ under the residue map,
\[
 \Pr(\widetilde E\mid K_\ell=k)
 \le\frac{\Pr(\widetilde E)}{\Pr(K_\ell=k)}.
\]
The unconditioned residue modulo $3^r$ has law $\mu_r$, by the same cap
reduction used in Theorem~\ref{thm:fixed-precision-ensemble}. This proves
\eqref{eq:ensemble-domination}. It is the asserted bound on the
essential likelihood ratio, and relative entropy is bounded by its
logarithm. The lattice local central limit theorem gives
$\Pr(K_\ell=k)\asymp_\delta\ell^{-1/2}$ in the central window.
\end{proof}

Theorem~\ref{thm:ensemble-divergence} applies at full precision, but its
direction is insufficient for condition $(?3)$. It proves that the
relative entropy per digit vanishes and rules out an exponential
likelihood separation from the canonical ensemble. O3 asks for a lower
bound on every residue, so the unresolved part is the lower tail of the
likelihood ratio.

For $r=o(\ell)$ and $\xi_0\in\mathbb Z/3^r\mathbb Z$, the theorem gives
\begin{equation}\label{eq:fixed-conductor-fourier}
 \widehat p_{\ell,k}(3^{\ell-r}\xi_0)
 =\widehat\mu_r(\xi_0)+o(1),
\end{equation}
uniformly in $\xi_0$ and in the central cost window.
Thus the fixed-cost law does not approach Haar at fixed precision.
Fixed-conductor modes may persist. For example, $\mu_1$ assigns masses
$1/3$ and $2/3$ to the two units modulo $3$, so its nontrivial Fourier
coefficient has modulus $1/\sqrt3$. Any Fourier estimate relevant to
O4 must separate these exact coarse modes and control frequencies whose
conductor is comparable with $\ell$.

\begin{empirical}[Fixed-precision projection test]
\label{thm:fixed-precision-finite}
We projected the exact fixed-cost tables onto $3^r$, for $r=1,2,3$,
and compared them in total variation with $\mu_r$ obtained from the
independent Syracuse recursion. At $\ell=12$ and $k=\ell$, the
distances were $0.02050$, $0.04024$, and $0.06228$. At
$k=\ell+\lfloor\sqrt\ell\rfloor$, they were $0.04310$, $0.07902$,
and $0.12677$. The distances decrease over the computed levels. This
finite test checks the implementation and is not used in the proof.
At full precision and $\ell=13$, the distances were $0.25195$ at
$k=\ell$ and $0.26078$ at
$k=\ell+\lfloor\sqrt\ell\rfloor$. The corresponding precision-one
values were $0.01885$ and $0.04003$. These full-precision values do not
determine an asymptotic limit.
\end{empirical}

\begin{empirical}[Finite microcanonical coverage]\label{thm:microcanonical-finite}
We computed the multiplicities $g_\ell(k,a)$ exactly through
$\ell=12$, checking at every level that their residue sum equals the
independent bounded-composition count. At $\ell=12$, the fractions of
unit residues reached at
$k=\ell-\lfloor\sqrt\ell\rfloor,\ell,
\ell+\lfloor\sqrt\ell\rfloor$ are respectively
$0.2735$, $0.8074$, and $0.9973$. The minimum multiplicity is still
zero in all three cases. A separate Boolean recursion extends the
support calculation through $\ell=16$. At that level the central and
upper-window coverage fractions are $0.9194$ and $0.99999972$, while
the first fixed cost covering every residue equals $\ell+5$ for each
$10\le\ell\le16$. These values neither prove nor refute the uniform
asymptotic lower bound.
\end{empirical}

\begin{empirical}[Fourier budget for fixed-cost multiplicities]
\label{thm:microcanonical-fourier}
For $p_{\ell,k}(a)=g_\ell(k,a)/\sum_b g_\ell(k,b)$, additive Fourier
inversion gives
\[
 \min_a\frac{p_{\ell,k}(a)}{1/(2\cdot3^{\ell-1})}
 \ge 1-\frac23\sum_\xi
 |\widehat p_{\ell,k}(\xi)-\widehat u_\ell(\xi)|.
\]
At the central cost $k=\ell$, the spectral sum grows from $12.11$ at
$\ell=4$ to $293.29$ at $\ell=12$ and is dominated by primitive
frequencies. At $k=\ell+\lfloor\sqrt\ell\rfloor$, it grows from $7.61$
to $122.63$. At $\ell=12$, the largest individual discrepancy is
$0.339$ at central cost and $0.280$ in the upper window, while the
largest primitive discrepancies are $0.0246$ and $0.00682$. Thus the
largest coefficients retain a coarse component even as each primitive
coefficient becomes smaller. The resulting lower bounds are negative
throughout. This
does not disprove uniform positivity; it shows that termwise absolute
summation of the finite Fourier series does not establish it on this
range.
\end{empirical}

Together, Empirical Results~\ref{thm:wcc} and~\ref{thm:conjecture3} and
Proposition~\ref{prop:beta-wcc} connect integer covering, the weighted
Syracuse distribution, and a real functional equation. The support
identity and the canonical decomposition
\eqref{eq:microcanonical} are exact. The required microcanonical
multiplicity estimate and Wirsching's Conjecture~2 remain unresolved.
All the finite-level computations in this subsection, from
Theorem~\ref{thm:microcanonical} through
Empirical Result~\ref{thm:microcanonical-fourier}, are in
\texttt{collatz-endogeny/\allowbreak sec9-wirsching-2003-conjecture3}.

\section{Two candidate lemmas, executed to completion}
\label{sec:lemmas}

Sections~\ref{sec:regimes}--\ref{sec:triangulation} identify two
natural next steps: a ``regime-2 lemma'' (decorrelation of subtree
aggregates at logarithmic precision) and a ``Z-number reduction lemma''
(formalizing the analogy with Littlewood--Offord/Bohr-set methods from
\cite{Tao2011}). The first yields an exact obstruction and a finite
test. The second retains an open frequency-stability step.

\subsection{The regime-2 lemma: structure theorem and a finite test of
its natural sufficient condition}

The originally sketched budget, $|\rho(\xi)|\le C_A D^{-A}$
\emph{uniformly in $\xi$}, contradicts Proposition~\ref{prop:tao123}
(an exact identity, not a lossy estimate): for any frequency $\xi$ of
primitive level $\ell'$ (i.e.\ $\xi=3^{\ell-\ell'}\xi_0$, $3\nmid\xi_0$)
and any depth $D\ge\ell'$, $\mathbb{E}[e(\xi F_D/3^\ell)] =
\varphi_{\ell'}(\xi_0)$ is \emph{independent of $D$}: depth beyond
$\ell'$ is a dummy parameter, because $3$-adic ultrametricity truncates
the relevant sum to the coarsest $\ell$ coordinates
(Proposition~\ref{prop:tao123}). A closed example at $\ell'=1$:
$\mathrm{Syrac}_1 = 2^{-a}\bmod 3 \in\{1,2\}$ with probabilities
$1/3,2/3$; $|\varphi(1)|=1/\sqrt3\approx 0.577$ for \emph{every} $D$,
with no decay whatsoever. The correct per-frequency budget is
$C_A\ell'^{-A}$ (Proposition~\ref{prop:tao117}, which requires
$3\nmid\xi$), not $C_AD^{-A}$; even with this correction, the total
$\ell^\infty$ budget summed over frequencies, $\sum_{\ell'\le \ell}
3^{\ell'}(C_A\ell'^{-A})^2 \asymp 3^\ell \ell^{-2A}$, never tends to
$0$: closing by frequency counting would require square-root
cancellation, $\ell^{-\ell'/2-\varepsilon}$ decay --- Regime~3
resurfacing inside what should have been the easy regime.

What survives is a genuine structure theorem. With $w_2 =
2^\Delta w_1 + c_\Delta$, $c_\Delta=(2^\Delta-1)/3$ (a unit multiplier
mod $3^\ell$) relating admissible sibling leaves at exponent gap
$\Delta$:

\begin{lemma}[Covariance spectral identity]\label{lem:cov-spectral}
For subtree-aggregate observables $Z_i:=N_{\eta_i}^{(m)}(w_i(v))$
(weighted populations truncated to precision $m$), with $v$ uniform on
admissible classes mod $3^{m+1}$,
\[
  \Cov(Z_1,Z_2) = 3^{-2m}\sum_{\xi\ne 0 \bmod 3^m}
  e(-\xi c_\Delta/3^m)\, S_1(-2^\Delta\xi)\, S_2(\xi),
\]
by character orthogonality and the affine sibling relation. The two
summation indices are independent by the tautology of expanding the
square (not an arithmetic hypothesis), but the dependence through the
common $v$ is not thereby cancelled: it survives entirely as the
diagonal frequency pairing $\xi_1=-2^\Delta\xi_2$.
\end{lemma}

\begin{proposition}[Exact coarse endogeny]\label{prop:exact-endogeny}
For complete populations ($\eta_i\equiv 1$) and precision $\ell=1$:
$\Corr(Z_1,Z_2)=+1$ if $\Delta\equiv 0\pmod 6$, and $=-1/2$ otherwise
--- independent of depth $D$ and of any truncation.
\end{proposition}

Proposition~\ref{prop:exact-endogeny} converts the endogeny barrier of
Theorem~\ref{thm:barrier} from a qualitative argument into a closed
form: aggregating leaves does \emph{not} decorrelate them; aggregation
projects onto the coarse $\sigma$-algebra generated by $w_i\bmod
3^\ell$, where the affine sibling relation of the fixed-pair regime
(\S\ref{sec:regimes}, Regime~1) reappears intact rather than vanishing.
This is a positive result with the opposite sign from what was hoped
for: what is provable is the persistence of coarse correlation, not its
disappearance.

Decomposing the sum of Lemma~\ref{lem:cov-spectral} by primitive level
and applying Cauchy--Schwarz plus Parseval gives $\Cov =
C_{\mathrm{exact}}(\ell_0;\Delta) + \mathrm{Err}$, with
$C_{\mathrm{exact}}$ exact, finite, and $D$-independent
(Proposition~\ref{prop:exact-endogeny}), and
\[
  \frac{|\mathrm{Err}|}{\mathbb{E}[Z_1]\mathbb{E}[Z_2]} \;\le\;
  K_m - K_{\ell_0}, \qquad K_\ell := 3^\ell \sum_y \mu_\ell(y)^2,
\]
the normalized \emph{collision mass} of the Syracuse measure $\mu$. A
conditional lemma would follow if $K_\infty := \lim_\ell K_\ell <
\infty$ (equivalently, the density $f=d\mu/d(\text{Haar})$ lies in
$L^2(\mathbb{Z}_3)$): the fine covariance, net of the exact coarse
component, would then be $\le K_\infty - K_{\ell_0}\to 0$, uniformly in
$m,D,\Delta$.

\begin{remark}[Consistency with the calibration of \S\ref{sec:calibration}]
\label{rem:calibration-consistency}
Proposition~\ref{prop:exact-endogeny} and the null result
$\rho_{\mathrm{eff}}\lesssim 0.06$ of \S\ref{sec:calibration} measure
different objects and do not conflict, though the point is worth making
explicit. The sharpest reason is structural, not merely quantitative:
in the synthetic model underlying \S\ref{sec:calibration}'s Arm~1 and
Arm~2, a child's type (its residue mod $3$ --- exactly the
$\ell=1$ object $Z_i$ of Proposition~\ref{prop:exact-endogeny}) is
computed as a deterministic function of the root's phase and the
branch index alone, reproducing the true cyclic admissibility
structure (validated against real integers), and is entirely
independent of the coupling intensity $\rho$ that Arm~2 varies: we
checked this directly in the simulator's code (2000 seeds,
$\rho\in\{0,0.3,0.7,1\}$, zero divergence in the resulting type
sequence), and confirmed that this deterministic type structure
reproduces the exact correlation of
Proposition~\ref{prop:exact-endogeny} ($+1$ for
$\Delta\equiv0\pmod6$, $-1/2$ otherwise) to four decimal places. In
other words, the coarse correlation is already present, identically,
in Arm~1 ($\rho=0$, the ``no coupling'' null) as it is in the real
tree; $\rho$ instead governs a disjoint axis --- whether sibling
subtrees' \emph{content} (the fresh randomness setting each subtree's
\emph{size}) is shared, not whether their root residues are
correlated. The bound $\rho_{\mathrm{eff}}\lesssim0.06$ therefore could
not have detected, or failed to detect,
Proposition~\ref{prop:exact-endogeny}'s correlation in the first
place: that correlation was never a form of coupling the calibration's
null hypothesis removes. A secondary, quantitative reason reinforces
this: even granting the coarse correlation full weight, its
contribution to the \emph{aggregate} variance the calibration measures
is diluted. Proposition~\ref{prop:exact-endogeny}'s covariance
$C_{\mathrm{exact}}(\ell_0;\Delta)$ is finite and $D$-independent,
while the aggregate means $\mathbb{E}[Z_i]$ grow exponentially with
truncation depth, so its normalized contribution
$C_{\mathrm{exact}}/(\mathbb{E}[Z_1]\mathbb{E}[Z_2])\to0$ as depth
increases --- consistent with the excess of \S\ref{sec:calibration}
decaying monotonically to zero with truncation depth. Separately, the
finite-level $L^2$ computation below
(Empirical Result~\ref{thm:l2-finite}) tests an upper-bound route,
not the fine covariance itself; \S\ref{sec:calibration} measures the
fine coupling directly and finds it small. Finally, the geometric branch
weights $q\cdot 2^{-a}$ concentrate on small gaps $\Delta$: weighting
pairs of branches by their natural second-moment mass
$2^{-a_i}\cdot2^{-a_j}$ (the same weighting entering $\Cov$ above),
$\Delta=2k$ has probability $3\cdot4^{-k}$, giving $\Pr[\Delta\equiv
0\!\!\pmod6]=1/21$ against $\Pr[\Delta\not\equiv0\!\!\pmod6]=20/21$
exactly, so the mean coarse correlation under this natural weighting is
exactly $\mathbb{E}[\Corr]=\tfrac{1}{21}(+1)+\tfrac{20}{21}(-\tfrac12)=
-\tfrac37$ (verified in
\texttt{collatz-endogeny/\allowbreak sec10-l2-refutation-and-jensen/}), consistent with the
observed negative sign of the residual excess in
\S\ref{sec:calibration}, though (per the structural point above) no
such quantitative match is actually required to reconcile the two
results. The two sections are therefore complementary rather than in
tension: Theorem~\ref{thm:barrier} and
Proposition~\ref{prop:exact-endogeny} establish that the recursion
\emph{cannot force} sibling decorrelation (the per-pair coarse
structure survives at every depth, in every model that respects the
true admissibility rule), while \S\ref{sec:calibration} establishes
that the actual aggregate residual, on the orthogonal axis of subtree
\emph{content}, is nonetheless empirically small.
\end{remark}

\begin{empirical}[Finite-level growth of the $L^2$ statistic]
\label{thm:l2-finite}
$K_\ell$ was computed exactly via the recursion
\begin{equation}\label{eq:mu-recursion}
  \mu_n(y) = \tfrac12\,\nu(2y\bmod 3^n) + \tfrac12\,\mu_n(2y\bmod 3^n),
\end{equation}
where $\nu$ is the law of $1+3F_{n-1}$ (derived below), solved as a
truncated geometric series along the cyclic orbit of multiplication by
$2$ mod $3^n$ (a single orbit, since $2$ is a primitive root mod $3^n$
for every $n$). For $\ell=1,\dots,17$, $K_\ell$ grows cleanly and
consistently \emph{linearly}, with increments converging monotonically
to $\approx 0.47$, with no sign of saturation ($K_1=5/3$ exactly,
matching a hand computation from $\Pr(\mathrm{Syrac}_1=1)=1/3$,
$\Pr(\mathrm{Syrac}_1=2)=2/3$). These finite data are evidence against
boundedness, but do not prove $K_\infty=\infty$.
\end{empirical}

\begin{proof}[Derivation of the recursion~\eqref{eq:mu-recursion}]
By Proposition~\ref{prop:tao123}, $F_n(a) =
2^{-a_1}\bigl(1+3F_{n-1}(a_2,\dots,a_n)\bigr)\bmod 3^n$. Write
$X:=1+3F_{n-1}(a_2,\dots,a_n)$, with law $\nu$, independent of $a_1$.
Conditioning on $a_1=1$ (probability $1/2$) gives $F_n = 2^{-1}X$;
conditioning on $a_1\ge 2$ (probability $1/2$), the memorylessness of
the geometric distribution gives $a_1-1\sim\mathrm{Geom}(2)$ again, so
$F_n = 2^{-1}\bigl(2^{-(a_1-1)}X\bigr)$, and the bracketed quantity has
exactly the law of $F_n$ itself (an independent copy). This gives the
self-referential fixed point $\mu_n(y) = \tfrac12\nu(2y) +
\tfrac12\mu_n(2y)$, all arithmetic mod $3^n$.
\end{proof}

We validated the recursion --- not merely its $\ell=1$ base case, which
would not exercise the recursive machinery at all --- against direct
Monte Carlo sampling of $\mathrm{Syrac}$ from its primary definition
\eqref{eq:syracuse-def} at $\ell=3,4$ ($3\times10^6$ samples each),
bin-by-bin: maximum discrepancy $0.000311$ ($\ell=3$) and $0.000213$
($\ell=4$), both within the expected Monte Carlo noise
($\approx 1/\sqrt{N}\approx 0.000577$).

\begin{empirical}[Finite-level $L^p$ collision spectrum]
\label{thm:lp-spectrum}
We evaluated \eqref{eq:lp-collision} for the Syracuse law at
$p\in\{1.1,1.25,1.5,1.75,1.9,2,2.1,2.5\}$ and $k\le14$, using the same
recursion and checks as Empirical Result~\ref{thm:l2-finite}. At $k=14$,
the values for $p=1.1,1.5,1.9,2,2.1,2.5$ were respectively
$1.129$, $2.186$, $5.771$, $7.743$, $10.603$, and $45.211$. Their last
consecutive ratios were $1.0017$, $1.0161$, $1.0512$, $1.0647$,
$1.0804$, and $1.1681$. Growth slows markedly below $p=2$ and increases
above it on this range. The data do not establish boundedness or
divergence for any $p$.
\end{empirical}

Both empirical results above, and the Monte Carlo validation of the
recursion, are in \texttt{collatz-endogeny/\allowbreak
sec10-l2-refutation-and-jensen}.

\begin{remark}[Connection to a conjectural tail exponent]
If the limiting Syracuse density exists and satisfies
$\Pr(f(U)>x)\sim Cx^{-2}$, then its $p$th moment is finite for $p<2$ and
infinite for $p\ge2$. The finite-level pattern in Empirical
Result~\ref{thm:lp-spectrum} is compatible with this threshold. It does
not prove the existence of the density, the tail asymptotic, or the
moment behavior at infinite level.
\end{remark}

\begin{remark}[Flatness questions for the Syracuse measure]
Tao's oscillation estimate, boundedness of $K_\infty$, and the
$L^\infty$-type conjecture $\beta=1$ are distinct statements. The
finite computation above does not settle the $L^2$ statement or either
of the other two.
\end{remark}

The regime-2 lemma is thus not an easier stepping stone toward
Regime~3: it is a sibling of it, sharing the same missing ingredient,
now with finite evidence against the natural sufficient condition.

\subsection{The Z-number reduction: a conditional dichotomy and an
exact negative result}

The second candidate lemma formalizes ``failure of the Weak Covering
Conjecture at scale $\ell$ with margin $c$ implies a Bohr-type
concentration configuration'', via the Littlewood--Offord/Riesz-product
technique of Tao's 2011 blog post \cite{Tao2011} (not a paper; we flag
this bibliographic correction explicitly, since it is a recurring
source of miscitation). That post proves conditional separation
results for $2^a-3^k$ via Halász's method \cite{Halasz1977} and inverse
Littlewood--Offord theory \cite{TaoVu2009,NguyenVu2011}, but
acknowledges that Baker's theorem already gives an unconditionally
stronger separation, so treats its own propositions as methodological
exercises. Mahler's classical ``Z-number'' problem \cite{Mahler1968}
concerns $\xi(3/2)^n$, a genuinely different (though structurally
adjacent) object from ours (which concerns $\times 2\bmod 3^s$ at
finite depth, not $\times 3/2\bmod 1$ on an infinite orbit); the honest
anchor for our object is the Erdős/Lagarias family, including the
general fractional-parts framework of
Flatto--Lagarias--Pollington
\cite{Erdos1979,Lagarias2009,FlattoLagariasPollington1995}.

\begin{definition}[Failure with margin $c$]
For $c>0$, the Weak Covering Conjecture \emph{fails at scale $\ell$
with margin $c$} if, for $j=\lceil(1+c)\log_4 3\cdot\ell\rceil$, some
$b\in(\mathbb{Z}/3^\ell\mathbb{Z})^\times$ lies outside the image of
$R_{j-1,j}$. Set $\gamma_c := (j+\ell)/\ell = 1+(1+c)\log_4 3$.
\end{definition}

\begin{definition}[Diagonal configuration]
A configuration of depth $s$ at slope $\gamma$, tolerance $\delta$, and
exception rate $\eta$ is $\xi\in\ZZ$, $3\nmid\xi$, such that in
$\ge(1-\eta)s$ of the scales $t\in\{1,\dots,s\}$ there is
$a\in[\gamma t - t^{2/3}, \gamma t+t^{2/3}]$ with
$\|\xi\cdot 2^a/3^t\|<\delta$.
\end{definition}

\begin{lemma}[Pre-wrap triviality]
For $\xi=1$, $\|2^a/3^s\|<\delta$ holds trivially for
$a\le\log_2 3\cdot s+\log_2\delta$ (the representative has not yet
wrapped around). Since $\gamma_c\ge 1.7925>\log_23\approx 1.585$ for
every $c\ge 0$, the forced configuration lives entirely in the
non-trivial post-wrap regime, by structural slack coming from the
$4^j$ vs.\ $3^\ell$ entropy accounting.
\end{lemma}

\begin{proposition}[Coarse modes and the sharp primitive hole barrier]
\label{prop:primitive-hole}
Let $N=3^\ell$, let $U=(\mathbb Z/N\mathbb Z)^\times$, and suppose
$\mu$ is any probability measure supported on $U$. Then
\[
 |\widehat\mu(N/3)|\ge\frac12.
\]
For $\ell\ge2$, fix $a\in U$ and let $u_a$ be uniform on
$U\setminus\{a\}$. This measure has a support hole, but for every
primitive additive frequency $3\nmid\xi$,
\[
 |\widehat u_a(\xi)|=\frac{1}{2\cdot3^{\ell-1}-1}.
\]
There are also probability measures with support holes for which every
primitive coefficient is zero.
\end{proposition}

\begin{proof}
At frequency $N/3$, the coefficient is
$p e(1/3)+(1-p)e(2/3)$, where $p$ is the mass in the residue class
$1\pmod3$. Its real part is $-1/2$, proving the first bound. For
primitive $\xi$, the Ramanujan sum over all units modulo $3^\ell$ is
zero when $\ell\ge2$. Removing $a$ leaves
\[
 \widehat u_a(\xi)=-\frac{e(\xi a/N)}{|U|-1},
\]
which proves the second assertion.
For the last assertion, take any probability measure on the units
modulo $3^{\ell-1}$ whose support is not full and lift it uniformly to
each of the three residues above every point. For $3\nmid\xi$, summing
the character over a fibre contributes
$1+e(\xi/3)+e(2\xi/3)=0$.
\end{proof}

\begin{proposition}[Primitive energy is fresh-fibre imbalance]
\label{prop:primitive-fibre-energy}
Let $\mu$ be a probability measure on $\mathbb Z/3^\ell\mathbb Z$.
For $b\bmod 3^{\ell-1}$ put
\[
 x_t(b)=\mu(b+t3^{\ell-1}),\qquad t=0,1,2.
\]
Then
\begin{equation}\label{eq:primitive-fibre-energy}
 \sum_{3\nmid\xi}|\widehat\mu(\xi)|^2
 =3^{\ell-1}\sum_b\bigl((x_0-x_1)^2+(x_1-x_2)^2+(x_2-x_0)^2\bigr).
\end{equation}
Thus all primitive coefficients vanish if and only if $\mu$ is the
uniform lift of its marginal modulo $3^{\ell-1}$. If
$\mu(x)=n_x/T$ for nonnegative integer counts and at least one fibre is
not uniform, then
\[
 \max_{3\nmid\xi}|\widehat\mu(\xi)|\ge\frac1T.
\]
For $\ell\ge2$, this bound is attained by the uniform law on all units
except one.
\end{proposition}

\begin{proof}
Parseval at level $\ell$ gives
\[
 \sum_\xi|\widehat\mu(\xi)|^2=3^\ell\sum_x\mu(x)^2.
\]
The nonprimitive coefficient at $\xi=3\eta$ is the coefficient at
$\eta$ of the marginal
$m(b)=x_0(b)+x_1(b)+x_2(b)$. Subtracting Parseval for that marginal
gives
\[
 3^{\ell-1}\sum_b
 \left(3\sum_{t=0}^2x_t(b)^2-m(b)^2\right),
\]
which is the right side of \eqref{eq:primitive-fibre-energy}. If an
integer fibre is nonuniform, the sum of its three squared pairwise
differences is at least two. There are $2\cdot3^{\ell-1}$ primitive
frequencies, so their maximum squared modulus is at least $T^{-2}$.
The final equality case follows from
Proposition~\ref{prop:primitive-hole}.
\end{proof}

The identity also shows that projective compatibility does not repair
the support-hole implication. Start with a non-full law at some coarse
level and lift it uniformly through every subsequent three-point
fibre. The resulting laws are projectively consistent and non-full,
but every new primitive spectrum is zero. Primitive concentration
measures imbalance of the newest digit, rather than inherited support
failure.

\begin{theorem}[Multiscale Parseval decomposition]
\label{thm:multiscale-parseval}
Let $(\mu_r)_{r\ge0}$ be compatible probability laws modulo $3^r$,
with $\mu_0$ the point mass on the trivial group. Define
\[
 K_r=3^r\sum_x\mu_r(x)^2,\qquad
 E_r=\sum_{3\nmid\xi}|\widehat\mu_r(\xi)|^2.
\]
Then
\begin{equation}\label{eq:multiscale-parseval}
 K_r-K_{r-1}=E_r,\qquad
 K_\ell=1+\sum_{r=1}^\ell E_r.
\end{equation}
If the laws are the projections of a probability measure $\mu$ on
$\mathbb Z_3$, then $\mu$ has an $L^2$ density with respect to Haar
measure if and only if
\[
 \sum_{r\ge1}E_r<\infty.
\]
\end{theorem}

\begin{proof}
Parseval identifies $K_r$ with the sum of
$|\widehat\mu_r(\xi)|^2$ over all characters modulo $3^r$. The
characters with $3\mid\xi$ are the pullbacks of the characters modulo
$3^{r-1}$. Compatibility identifies their coefficients with those of
$\mu_{r-1}$, so their total contribution is $K_{r-1}$. The remaining
characters are primitive and contribute $E_r$. This proves
\eqref{eq:multiscale-parseval} by subtraction and iteration.

The level-$r$ density relative to Haar measure is the conditional
density martingale of $\mu$, and its squared $L^2$ norm is $K_r$.
Boundedness of these norms is equivalent to $L^2$ martingale
convergence. The limit is an $L^2$ density for $\mu$. Conversely, if
$\mu$ has an $L^2$ density, the level densities are its conditional
expectations and have bounded $L^2$ norms. Equation
\eqref{eq:multiscale-parseval} completes the equivalence.
\end{proof}

Theorem~\ref{thm:multiscale-parseval} identifies the precise overlap
between O5 and the $L^2$ branch of O7. Proposition
\ref{prop:primitive-fibre-energy} describes the increment created by
the newest digit, while $K_\ell$ is the accumulated energy through
level $\ell$. The identity neither selects the largest primitive
coefficient at one level nor aligns frequencies across levels, so it
does not prove Conjecture~\ref{lem:B}.

For the $q$-adic Syracuse laws at the geometric weight $\alpha=1$, the
one-level collision admits a second diagonalization. It uses Fourier
analysis after ordering the preceding law by a discrete logarithm,
rather than additive Fourier analysis on the residue group.

\begin{theorem}[$q$-adic logarithmic collision renormalization]
\label{thm:logarithmic-collision}
Let $q$ be an odd prime, put $d=\ord_q(2)$, and suppose that
$v_q(2^d-1)=1$. Let $\mu_{q,\ell}$ be the Syracuse law modulo
$q^\ell$ with digit probabilities $\Pr(a)=2^{-a}$, and define
\[
 K_{q,\ell}=q^\ell\sum_x\mu_{q,\ell}(x)^2.
\]
For $M=q^{\ell-1}$, put
\[
 a_j=\mu_{q,\ell-1}\left(\frac{2^{dj}-1}{q}\bmod M\right),\qquad
 A_m=\sum_{j=0}^{M-1}a_j e^{-2\pi i m j/M}.
\]
Then, for every $\ell\ge1$,
\begin{equation}\label{eq:logarithmic-collision}
 K_{q,\ell}=\sum_{m=0}^{M-1}
 \frac{q(4^d-1)|A_m|^2}
 {3(4^d+1-2^{d+1}\cos(2\pi m/M))}.
\end{equation}
In particular, with
\[
 c_{q,d}=\frac{q(2^d-1)}{3(2^d+1)},\qquad
 C_{q,d}=\frac{q(2^d+1)}{3(2^d-1)},
\]
one has
\begin{equation}\label{eq:logarithmic-bounds}
 c_{q,d}K_{q,\ell-1}\le K_{q,\ell}
 \le C_{q,d}K_{q,\ell-1}.
\end{equation}
For $q\ge5$, $c_{q,d}\ge35/27>1$. Hence
$K_{q,\ell}\ge c_{q,d}^{\ell}$, and the projective Syracuse law has
no $L^2$ density on $\mathbb Z_q$.
\end{theorem}

\begin{proof}
The lifting assumption gives $\ord_{q^\ell}(2)=dM$. Write
$y_k=2^k\bmod q^\ell$, $x_k=\mu_{q,\ell}(y_k)$, and let
$n_k$ be the mass at $y_k$ of the law of $1+qF_{\ell-1}$. The Syracuse
recursion gives
\[
 x_k=\frac12 n_{k+1}+\frac12x_{k+1},
\]
with indices modulo $dM$. Also, $n_{dj}=a_j$, and $n_k=0$ when
$d\nmid k$. The map $j\mapsto(2^{dj}-1)/q\bmod M$ is a permutation:
equality of two outputs modulo $M$ is equivalent to equality of the
corresponding powers modulo $q^\ell$, where $2^d$ has order $M$.

Set $u_j=x_{dj}$. A block of $d$ recursion steps gives
\[
 u_j=2^{-d}(a_{j+1}+u_{j+1}),\qquad
 x_{dj+r}=2^r u_j\quad(0\le r<d).
\]
Hence $\sum_kx_k^2=(4^d-1)\sum_ju_j^2/3$. At frequency
$z=e^{2\pi i m/M}$, the circular filter from $a$ to $u$ has multiplier
$z/(2^d-z)$, whose squared modulus is
$(4^d+1-2^{d+1}\cos(2\pi m/M))^{-1}$. Parseval and $q^\ell=qM$ prove
\eqref{eq:logarithmic-collision}. Another application of Parseval gives
$K_{q,\ell-1}=\sum_m|A_m|^2$. The minimum and maximum of the multiplier
occur at cosine $-1$ and $1$, respectively, which proves
\eqref{eq:logarithmic-bounds}. If $q\ge5$, then $d\ge3$, so
$c_{q,d}\ge(5/3)(7/9)=35/27$.
\end{proof}

For $q=3$, one has $d=2$, and \eqref{eq:logarithmic-collision} reduces
to
\[
 K_\ell=\sum_m\frac{15}{17-8\cos(2\pi m/M)}|A_m|^2.
\]
Together with Theorem~\ref{thm:multiscale-parseval}, this gives
\[
 E_\ell=\sum_m
 \left(\frac{15}{17-8\cos(2\pi m/M)}-1\right)|A_m|^2.
\]
The signed multiplier is positive exactly when
$\cos(2\pi m/M)>1/4$. Thus the unresolved $q=3$ part of O7 is a
balance between low and high frequencies of the preceding law in
logarithmic order. The uniform lower bound is only $3/5$, so the
exponential argument for $q\ge5$ does not apply.

Proposition~\ref{prop:primitive-hole},
Proposition~\ref{prop:primitive-fibre-energy},
Theorem~\ref{thm:multiscale-parseval}, and
Theorem~\ref{thm:logarithmic-collision} above are all verified in
\texttt{collatz-endogeny/\allowbreak sec10-primitive-spectral-analysis},
including the general-$q$ case for
$q=7,11,13,17,19$.

\begin{definition}[Primitive spectral concentration
$SC_{\rm prim}(\varepsilon)$]
$\mu_{\ell,j}$ (the law of $Y=\sum R\bmod 3^\ell$) satisfies
$SC_{\rm prim}(\varepsilon)$ if there is $\xi$ with $3\nmid\xi$ and
$|\hat\mu_{\ell,j}(\xi)|\ge 3^{-\varepsilon\ell}$.
\end{definition}

Allowing every nonzero frequency would make the condition automatic by
Proposition~\ref{prop:primitive-hole}. The primitive restriction is also
required by the diagonal configuration. The same proposition shows
that a support hole alone cannot imply $SC_{\rm prim}(\varepsilon)$ for
any fixed $\varepsilon<1$, or even force a nonzero primitive
coefficient.

\begin{definition}[Chained configuration]
A chained configuration of length $L$ has a primitive frequency $\xi$,
an interval $I$ of $L$ consecutive scales, and exponents $a_t$ for
$t\in I$ satisfying
\[
 |a_t-\gamma t|\le t^{2/3},\qquad
 \|\xi2^{a_t}/3^t\|<\delta,\qquad
 |a_{t+1}-a_t|\le C.
\]
The constants $\delta<1/2$ and $C$ do not depend on $L$.
\end{definition}

\begin{conjecture}[Primitive chained concentration target]
\label{lem:B}
For fixed $c>0$ and sufficiently small $\varepsilon>0$, failure at
scale $(\ell,c)$ together with $SC_{\rm prim}(\varepsilon)$ forces a
chained configuration at slope $\gamma_c$ whose length
$L=L_\ell$ satisfies $L/\log\ell\to\infty$.
\end{conjecture}

The unchained diagonal condition above does not control how the
exponents at adjacent scales fit together and is insufficient for this
target. The Riesz-product and Halász route supplies good bases on a set
whose mass may be exponentially small when the Fourier coefficient is
$3^{-\varepsilon\ell}$; with fixed $\varepsilon$, the available
intra-scale argument yields only bounded chains. The alternative
recursion in $j$ is a convex Pascal mixture and does not keep a
concentrating frequency fixed. Neither route proves
Conjecture~\ref{lem:B}, which is not used elsewhere in the paper.
Proposition~\ref{prop:primitive-fibre-energy} supplies an exact
intermediate quantity: a proof from support data must first produce
quantitative fresh-fibre imbalance at consecutive scales, and must then
align the resulting primitive frequencies into one chain.

\subsection{Proposition C: an exact wall of constants}

\begin{proposition}[Unrestricted Halász deficit]\label{prop:halasz-deficit}
The Fourier-inversion argument for a support hole gives a coefficient
of scale at least $3^{-\ell}$ at some nonzero frequency, but that
frequency may be coarse and cannot be used in the diagonal
configuration. No primitive lower bound follows from a hole, by
Proposition~\ref{prop:primitive-hole}. Even if a primitive coefficient
of size $3^{-\ell}$ is added as a hypothesis, Halász's theorem yields
the budget $\sum\|\cdot\|^2\le 0.549\ell$
against a required ceiling of $0.25\ell$ ($\delta=1/2$, $d=\ell$
generators): a deficit of a factor $\ge 2.2$, rising to $\approx 7$ at
a realistic generator density $\rho\approx 0.3$. In Tao's original 2011
application this closes because $\log q\ll d$ (unbounded ratio); in
the Weak Covering Conjecture regime, $\log_2 q/d\ge 1.6$ --- the
technique does not merely degrade, it inverts sign.
\end{proposition}

\begin{theorem}[Jensen identity and the exact Fourier-budget deficit]
\label{thm:jensen}
Let $Z=\sum_{g\ge1} w_g\, e(U_g)$, $U_g$ i.i.d.\ uniform, $w_g =
p(1-p)^{g-1}$ ($\sum_g w_g=1$), with $p\ge\tfrac12$. Then, exactly,
\begin{equation}\label{eq:jensen}
  \Lambda := \mathbb{E}\bigl[\log(1/|Z|)\bigr] = \log(1/p).
\end{equation}
\end{theorem}

\begin{proof}
By the memoryless self-similarity of the geometric weights
($w_g/(1-w_1)=p(1-p)^{g-2}$ for $g\ge2$, the same geometric sequence
re-indexed), $Z \overset{d}{=} p\, e(U_1) + (1-p)\, Z'$, with $Z'$ an
independent copy of $Z$; since $|Z'|\le 1$ always (a convex combination
of unimodular terms), and $p/(1-p)\ge 1$ for $p\ge1/2$, we have
$|Z'|\le p/(1-p)$ whenever $p\ge1/2$. Writing $c:=(1-p)Z'/p$, so
$|c|\le1$, factor $p\,e(U_1)+(1-p)Z' = p\,e(U_1)\bigl(1+c\,e(-U_1)\bigr)$
and take $\log|\cdot|$: the function $z\mapsto\log|1+cz|$ is (the real
part of) a holomorphic function of $z$ with no zero inside the disc
$|z|<1/|c|\ni$ the unit circle (as $|c|\le1$), so by the mean-value
property (Jensen's formula with no zeros enclosed) its average over
$z=e(U_1)$, $U_1$ uniform, equals its value at $z=0$, namely $0$.
Hence $\mathbb{E}_{U_1}\log|p\,e(U_1)+(1-p)Z'| = \log p$ for every
realization of $Z'$, and averaging over $Z'$ gives
$\mathbb{E}[\log|Z|]=\log p$, i.e.\ \eqref{eq:jensen}.
\end{proof}

Applying Theorem~\ref{thm:jensen} with $p=p_c=1/\gamma_c\approx0.558$
(the tilted model matching the Weak Covering Conjecture's threshold
slope; note $p_c\ge\tfrac12$, as the theorem requires) gives
$\Lambda = \log\gamma_c \approx 0.5836+O(c)$, confirmed by Monte Carlo
($8\times10^6$ samples: $0.58344\pm0.0005$), against the value
$\log 3\approx 1.0986$ that would be needed to exclude a diffuse hole
by an annealed $\ell^1$ Fourier-budget argument: a deficit factor of
$\approx 1.88$. The naive threshold for inversion, $p=1/3$, sits below
$\tfrac12$, where Theorem~\ref{thm:jensen}'s identity no longer holds:
the same stationary-phase argument gives $\Lambda(p)<\log(1/p)$ strictly
for $p<\tfrac12$, so the actual inversion point is $\gamma\approx3.31$
($\Lambda\approx1.032$ at $p=1/3$, against $\log3\approx1.0986$), not
$\gamma=3$, still far above the regime of the conjecture
($j^\ast(\ell)/\ell\approx1.2$, Empirical Result~\ref{thm:wcc}). No
member of the wider $\ell^r$ family ($r\ge1$) closes the gap either:
the closing criterion is the same $\|Z\|_r<1/3$ for every $r$, and
$\|Z\|_r$ increases with $r$, so $r=1$ is already the best choice and
$r\to0$ (the Jensen exponent above) is a strict lower bound for the
whole family (verified in
\texttt{collatz-endogeny/\allowbreak sec10-l2-refutation-and-jensen/}).

\begin{proposition}[Annealed budget for a diffuse hole]
\label{thm:propC}
A primitively spectrally diffuse failure of the Weak Covering Conjecture is
compatible with the specific annealed $\ell^1$ Fourier budget above,
with slack $\approx 1.88$. That budget therefore does not exclude such
a failure. This branch is the one
identified alongside $\beta=1$ (\S\ref{sec:triangulation}) and the
finite-level $L^2$ growth (Empirical Result~\ref{thm:l2-finite}): a
further articulation, within the same second-moment Fourier vocabulary
as the $L^2$ condition. Theorem~\ref{thm:multiscale-parseval} identifies
their exact telescoping relation, while both enter the covariance sum
$\Cov\propto\sum_\xi S_1(\xi)S_2(\xi)^\ast$ of
\S\ref{sec:regimes}), of the same missing ingredient --- not a
methodologically independent fourth or fifth route in its own right.
\end{proposition}

We emphasize the correct framing, consistent with the fragilities
listed above: this is a reduction of the spectrally \emph{concentrated}
branch to the Erdős--Lagarias--Furstenberg family, conditional on the
open stability step, together with a calculation showing that one
annealed $\ell^1$ budget does not exclude the diffuse branch. It is not
a reduction of the Weak
Covering Conjecture to the classical Z-number problem, which
Theorem~\ref{thm:propC} and the discussion above show would be
imprecise.

\section{Contextualizing the barrier in the wider literature}
\label{sec:context}

A directed literature search, going beyond the Collatz-specific
literature into the general theory of Fourier decay for self-similar
measures, situates the barrier precisely and rules out several
plausible-looking shortcuts, each for a specific, verified reason.

\subsection{The Syracuse measure is a rigid $3$-adic self-similar
measure}

Unrolling the renewal recursion of \eqref{eq:syracuse-def} gives, in
the limit,
\begin{equation}\label{eq:X-selfsimilar}
  X = \sum_{k\ge1} 3^{k-1}\, 2^{-S_k}, \qquad S_k:=a_1+\cdots+a_k,
\end{equation}
the fixed point of the countable affine iterated function system
$\phi_a(x) = 2^{-a}(1+3x)$ on $\mathbb{Z}_3$, weighted by
$\Pr(a)=2^{-a}$, every map sharing the \emph{same} ultrametric
contraction ratio $|3\cdot2^{-a}|_3 = 1/3$. This is a genuine (if
non-standard, countable, overlapping) self-similar structure, so the
question of whether the general theory of Fourier decay for
self-similar measures \cite{BakerBanaji2026,Bremont,LiSahlsten2022}
transfers is a fair one to ask, and we checked it explicitly.

\begin{proposition}[Rigidity of $\mathbb{Z}_3^\times$]\label{prop:rigidity}
Every infinite closed subgroup of $(\mathbb{Z}_3,+)$ has finite index
(it equals $3^k\mathbb{Z}_3$ for some $k\ge 0$); consequently, since
$\mathbb{Z}_3^\times\cong\mathbb{Z}/2\times\mathbb{Z}_3$, every
infinite closed subgroup of $\mathbb{Z}_3^\times$ has finite index.
Since $2$ is a topological generator of $\mathbb{Z}_3^\times$ (a
primitive root modulo $3^n$ for every $n$), the unit phases
$2^{-a}$ spread maximally, the opposite of the ``thin multiplicative
orbit'' phenomenon (a Pisot-type obstruction to the Rajchman property)
that drives non-uniform Fourier decay on the real line.
\end{proposition}

\begin{theorem}[The general theory does not transfer]\label{thm:no-transfer}
\begin{enumerate}[label=(\alph*)]
\item The pushforward mechanism of Baker--Banaji
  \cite{BakerBanaji2026} (Archimedean stationary phase, requiring a
  $C^2$ map with $f''\ne 0$ off a null set) would give, at best,
  \emph{qualitative} Rajchman decay for the analogue considered here
  --- strictly weaker than the super-polynomial decay already
  guaranteed by Proposition~\ref{prop:tao117}.
\item Brémont's Pisot-type criterion for non-Rajchman self-similar
  measures, transported to $\mathbb{Z}_3$, trivializes by
  Proposition~\ref{prop:rigidity}: there is no room for a
  ``thin'' $3$-adic multiplicative orbit of the kind the criterion
  requires.
\item The renewal-theoretic technique of Li--Sahlsten
  \cite{LiSahlsten2022}, which needs \emph{two distinct} contraction
  ratios with a Diophantine relation between their logarithms, also
  degenerates: every branch of \eqref{eq:X-selfsimilar} has the
  \emph{same} $3$-adic contraction ratio $1/3$, so the relevant
  ``scale walk'' is deterministic ($S_n = n\log 3$, a pure lattice
  walk with no overshoot fluctuation), not the non-lattice random walk
  the technique is built to exploit.
\end{enumerate}
\end{theorem}

We stress that Baker--Banaji's general phenomenon --- Rajchman
measures with no uniform decay rate are generic, not pathological, in
the class of self-similar/self-conformal measures --- remains a fair
piece of context: it shows that the kind of failure exhibited by
Theorem~\ref{thm:propC} has real structural precedent in the reference
literature, legitimizing the framing that a proof of $\beta=1$/the Weak
Covering Conjecture genuinely needs the specific arithmetic input of
this problem, not generic smoothness. It should not, however, be
over-interpreted as evidence \emph{for} the diffuse scenario itself:
Baker--Banaji's construction requires fine-tuning of a continuous
parameter (a Liouville-type construction, dense in the family but
exceptional in measure; typical parameters give polynomial decay
\cite{Solomyak1995,VarjuYu2022}), while the Syracuse measure is a rigid
arithmetic object with no tunable parameter, and
Proposition~\ref{prop:rigidity} shows the enabling mechanism (fine
Diophantine approximation) simply does not exist on the $3$-adic side
--- which cuts against the diffuse scenario as much as for it. We cite
Theorem~\ref{thm:no-transfer} with this disanalogy stated explicitly.

\subsection{A related fixed-modulus proposal}

Chang \cite{Chang2026} studies the compressed odd-to-odd orbit through
$X_t=\mathbb{1}[n_t\equiv1\pmod4]$ and the last index of each maximal
run of $X_t=1$. Its Map Balance Theorem gives a map-level discrepancy
of one between two gap classes. The remaining proposed condition asks
for a bound, with an unspecified threshold $\delta_{\max}$, on the
imbalance between residues $9$ and $25$ modulo $32$ at those times.

The claimed implication to Collatz is not self-contained in
\cite{Chang2026}: the threshold is delegated to a block-TV budget in a
companion work. The current version of that companion work
\cite{ChangCompanion2026} instead gives a conditional implication from
equidistribution at growing moduli, together with uniform control of
unbounded gap tails. Fixed-modulus balance does not supply those stated
hypotheses. We therefore use the modulo-$32$ statement only as a related
empirical diagnostic, not as a reduction or as an open problem of this
paper.

We tested this condition directly on real orbits. An ensemble of many
moderate orbits ($2$--$5\times10^4$ roots, $20$--$50$ bits) shows the
aggregate deviation plateauing at $1.2$--$1.9\%$, which is on its own
consistent with the conjecture (only boundedness is claimed, not
convergence to zero); a more faithful test --- a single very long
orbit (a random $16000$-bit start, $38395$ compressed steps, $4761$
relevant burst-endings) --- shows the deviation tracking the
$1/\sqrt{i}$ statistical-noise curve closely, with no systematic trend,
ending at a deviation of $0.00053$. Five further orbits from $500$ to
$8000$ bits show the same pattern. These measurements concern only the
fixed-modulus statistic. They do not test growing-modulus
equidistribution or uniform integrability of gap lengths.

\section{Discussion: why a precisely characterized barrier is a
genuine contribution}
\label{sec:discussion}

A natural objection to the material of
\S\S\ref{sec:endogeny}--\ref{sec:context} is that neither candidate
lemma of \S\ref{sec:lemmas} produced a proof, so the paper might seem
to document an absence rather than a result. We think this objection,
while natural, does not survive scrutiny of what was actually
produced.

First, the outcome of executing both candidate lemmas was not
silence: it was two independent, rigorous negative/structural results,
each of which pins down \emph{precisely} why a natural approach fails
--- Proposition~\ref{prop:exact-endogeny}'s exact closed-form
persistence of coarse correlation; Empirical
Result~\ref{thm:l2-finite}'s finite-level test of a specific
sufficient condition;
Theorem~\ref{thm:jensen}'s exact identity quantifying, to a specific
constant factor ($\approx1.88$), why a whole family of Fourier-analytic
techniques cannot close the gap. Precisely characterizing an
obstruction --- showing not merely that an approach fails, but
\emph{exactly} why, with a quantified margin --- is itself a
recognized and citable form of mathematical contribution, independent
of whether the underlying conjecture is eventually resolved.

Second, these results do not stand alone: a related underlying
arithmetic statement --- some form of equidistribution or
near-independence of $3$-adic digits along the true, deterministic
integer dynamics, not merely along an idealized i.i.d.\ model --- is
the crux reached by four complementary vocabularies:
(1) the smoothing-transform/endogeny route (\S\ref{sec:endogeny}); (2)
the support identity between Wirsching's Weak Covering Conjecture and
Tao's weighted $\beta=1$ problem in Proposition~\ref{prop:beta-wcc};
(3) Wirsching's 2003 real functional
equation and its Fabius density (\S\ref{sec:triangulation}); (4) the
second-moment Fourier analysis underlying both the $L^2$ condition and
the spectral dichotomy of Proposition~C (\S\ref{sec:lemmas}), which
share the common character decomposition
$\Cov\propto\sum_\xi S_1(\xi)S_2(\xi)^\ast$ and the exact
multiscale identity of Theorem~\ref{thm:multiscale-parseval}. The WCC
and $\beta=1$ formulations share an exact support and cost identity,
but differ by geometric multiplicity. The $L^2$ and spectral-dichotomy
branches share the same second-moment character decomposition. A
barrier visible from only one angle is a candidate
artifact of that technique; a barrier reached by four complementary
vocabularies (branching-process fixed points, integer
covering/entropy counting, a real functional equation, second-moment
character sums, and elementary residue counting) is evidence that the obstruction
is a structural feature of the problem itself.

Third, this theoretical package sits alongside a result
(Empirical Result~\ref{thm:kl}) whose validity does not depend on any of the
above framing at all: a direct, falsifiable, statistically calibrated
finite-window measurement comparing two predictions in the
literature. We consider this the single most framing-independent
contribution of the paper, and recommend it be weighted accordingly
by a reader assessing the paper's overall reliability.

We therefore present this paper's contribution not as an attempted
proof of any part of the Collatz conjecture, nor as a report of
failure, but as a precise map of where a specific, natural,
technically rich line of attack --- Tao's Syracuse-measure program,
generalized to $qx+1$ --- runs into a wall, together with an exact
quantification of the size and shape of that wall from several
angles, and one exact finite-window empirical result
obtained along the way.

\section{Conclusion and open directions}
\label{sec:conclusion}

We summarize the eight tracked points. O6 is an audit closure; the
other seven contain open mathematical statements.

\begin{enumerate}[label=(O\arabic*)]
\item \textbf{Cancellation of centered fine-frequency correlations
  across sibling subtree aggregates} (Theorem~\ref{thm:barrier},
  Remark~\ref{rem:barrier-reading}). Pairwise fresh-digit independence
  is impossible: Theorem~\ref{thm:fresh-digit-coupling} gives total
  variation $1-3^{-s}$ and mutual information $s\log3$ for every fixed
  path pair. The open arithmetic statement concerns cancellation only
  after independent averaging over the two path indices and removal of
  the exact coarse affine modes.
\item \textbf{Wirsching's Weak Covering Conjecture} (Conjecture~\ref{conj:wcc}),
  and the stronger weighted covering estimate needed for Tao's
  $\beta=1$ (\S\ref{sec:triangulation}). Proposition~\ref{prop:beta-wcc}
  identifies their common support. Theorem~\ref{thm:wcc-large-deviation}
  proves that even ideal multiplicity inside the WCC cost slice remains
  exponentially insufficient; the weighted estimate must reach the
  typical cost window.
\item \textbf{Wirsching's Conjecture 2} (\S\ref{sec:triangulation}),
  above the numerically tested Conjecture~3. Conjecture~1 is proved in
  Theorem~\ref{thm:wirsching-conj1}. Theorem~\ref{thm:microcanonical}
  proves that its target condition $(?3)$ implies the weighted
  cylinder estimate in O2. Proposition~\ref{prop:complex-deconditioning}
  identifies the conditional local limit theorem needed to pass from
  canonical weights to fixed cost. Theorem~\ref{thm:ensemble-divergence}
  shows that their likelihood gap is at most polynomial at every
  precision, but in the direction opposite to the required lower bound.
\item \textbf{Exponential control of the fine-frequency tail at depth
  $\ell\asymp D$} (\S\ref{sec:regimes}, Regime~3), after separating
  sublinear-conductor modes. Theorem~\ref{thm:sublinear-precision-ensemble}
  proves that those modes are approximated by the corresponding
  coefficients of $\mu_r$. Fixed-conductor coefficients need not
  decay. Theorem~\ref{thm:linear-block-nonequivalence} shows that
  equivalence already fails at linear scale for the underlying cost
  vectors, although this separation need not survive the residue map.
  The remaining linear-precision estimate is the input
  required by the second-moment route; Remark~\ref{thm:regime3}
  records nearby open theories without claiming an equivalence.
\item \textbf{Exclusion of primitively spectrally diffuse Weak Covering
  Conjecture failures}. Proposition~\ref{prop:primitive-hole} proves
  that a support hole alone need not produce any primitive Fourier
  coefficient. Proposition~\ref{prop:primitive-fibre-energy}
  identifies primitive Fourier energy exactly with imbalance among the
  three fresh children of each residue. Projective compatibility and
  inherited holes still allow zero primitive spectrum.
  Proposition~\ref{thm:propC} shows that the
  annealed $\ell^1$ budget tested here has insufficient strength; it
  does not rule out other $\ell^1$ arguments.
\item \textbf{Fixed-modulus residue balance, audit closed}
  (\S\ref{sec:context}). The statistic proposed in \cite{Chang2026}
  was tested here, but its claimed implication depends on an unspecified
  threshold and on a companion framework whose current sufficient
  hypothesis uses growing moduli and tail control
  \cite{ChangCompanion2026}. It is not counted as an open problem of
  this paper.
\item \textbf{Absolute continuity and the tail index of the tilted
  Syracuse density} in the coherent $q$-adic root environment
  (Conjecture~\ref{conj:tail-index}), for every odd $q\ge3$.
  Theorem~\ref{thm:lp-collision} characterizes its $L^p$ form by
  normalized collision probabilities. At $p=2$,
  Theorem~\ref{thm:multiscale-parseval} identifies bounded collision
  mass with summability of the primitive energies from O5. The matching
  identity in Theorem~\ref{thm:logarithmic-collision} gives exponential
  collision growth for the nonexceptional primes $q\ge5$ and expresses
  the remaining $q=3$ increment as an explicit low-versus-high spectral
  balance. It does not yet give a uniform sign margin at $q=3$. The matching
  i.i.d. statement is
  Theorem~\ref{thm:iid-tail}; its arithmetic transfer remains open.
  Empirical
  support, however, has substantially strengthened since the original
  under-powered measurement: a $20\times$ larger sample ($10^5$ roots,
  the pre-registered comparison called for above) produced, for the
  first time in this investigation, a clean GPD threshold-stability
  plateau at the predicted value across every threshold level tested,
  a tightly concentrated bias-corrected Hill estimate consistent with
  it across all headroom levels, and a Vuong test against a truncated
  lognormal alternative that no longer favors the alternative. Two
  calibration checks (against an exact synthetic Pareto sample of the
  predicted index, and against a deliberately mis-specified
  normalization exponent) found no artifact behind this shift. We
  regard the conjecture as now strongly supported for $q=5$, though
  not proved: the Vuong test yields non-rejection rather than
  confirmation, and the finite-depth factors remain pre-asymptotic at
  the headroom reached (their median continues to
  drift with headroom even as the estimated tail index does not).
\item \textbf{The exponent transition on the genuine arithmetic tree}
  (Conjecture~\ref{conj:transition-arithmetic}): established here only
  for the matching i.i.d.\ branching-random-walk model
  (Theorem~\ref{thm:transition-model}, and, conditionally,
  Proposition~\ref{prop:transition-fine}); for $q=3$ this is the
  Growth Exponent Conjecture of Kontorovich--Lagarias and
  Applegate--Lagarias, open since 1995, of
  which the best unconditional partial result remains the one-sided
  bound of Krasikov--Lagarias \cite{KrasikovLagarias2003}.
\end{enumerate}

Any progress on any one of (O1)--(O8) would very likely transfer, given
the structural identities established in this paper, to progress on
several of the others simultaneously. We regard this convergence as
the paper's central finding.

\section*{Data Availability Statement}

Public verification scripts for every numerical claim in this paper,
organized by section, are at
\url{https://github.com/faculdade/collatz-endogeny}. This includes the
pressure and freezing computations of \S\ref{sec:pressure}, the
residual-coupling experiment of \S\ref{sec:calibration}, the
Kontorovich--Lagarias vs.\ Volkov gate of \S\ref{sec:kl}, the Weak
Covering Conjecture test of \S\ref{sec:wcc}, the Wirsching
Conjecture~3 test, the $K_\ell$/Jensen computations, the Chang
residue-balance test, and the checks for
Theorem~\ref{thm:wcc-large-deviation},
Theorems~\ref{thm:fresh-digit-coupling} and
\ref{thm:linear-block-nonequivalence},
Propositions~\ref{prop:primitive-hole} and
\ref{prop:primitive-fibre-energy},
Theorems~\ref{thm:multiscale-parseval} and
\ref{thm:logarithmic-collision}, and Empirical
Results~\ref{thm:lp-spectrum}, \ref{thm:microcanonical-finite},
\ref{thm:microcanonical-fourier}, and \ref{thm:fixed-precision-finite}.

\section*{Conflict of Interest}

The author declares that there is no conflict of interest regarding
the publication of this paper.

\section*{Acknowledgments}

This research was conducted with the assistance of AI systems, used
to support literature review, mathematical analysis, numerical
experimentation, and manuscript preparation, following what we
believe is best practice in the small existing literature on
AI-assisted contributions to this specific problem.

\appendix

\section{Numerical validation methodology}
\label{app:validation}

We record, for completeness, the validation discipline applied to each
computational result reported above, since an unvalidated recursive
computation is a known source of silent error.

\begin{itemize}
\item \emph{Pressure equation roots} (\S\ref{sec:pressure}): solved by
  bisection independently of the closed-form derivation; cross-checked
  against empirical counting slopes on real reverse trees for $q=5,7$.
\item \emph{Weak Covering Conjecture DP} (\S\ref{sec:wcc}): validated
  against an independently computed reference table for
  $\ell=1,\dots,7$ and a hand-checked point ($\ell=2,j=2$).
\item \emph{WCC large-deviation barrier}
  (Theorem~\ref{thm:wcc-large-deviation}): the exact negative-binomial
  tail was evaluated in log space through $\ell=1000$ and compared
  with the closed-form rate $I(1+\log_4 3)$. The computation is a
  check on the constants; the proof uses the exact mass formula and
  Stirling's formula.
\item \emph{Primitive Fourier hole barrier}
  (Proposition~\ref{prop:primitive-hole}): direct character sums through
  $\ell=8$ reproduce $1/(2\cdot3^{\ell-1}-1)$ for every primitive
  coefficient of the one-hole unit law. They also give zero, up to
  floating-point error below $6\times10^{-13}$, for every primitive
  coefficient of a non-full law lifted uniformly from the preceding
  modulus.
\item \emph{Primitive fibre energy}
  (Proposition~\ref{prop:primitive-fibre-energy}): random integer count
  vectors through $\ell=8$ satisfy \eqref{eq:primitive-fibre-energy}
  with absolute discrepancy below $3\times10^{-16}$. Uniform lifts
  give zero primitive energy at every tested level. The computation
  also checks the sharp integer-count bound $1/T$.
\item \emph{Multiscale Parseval decomposition}
  (Theorem~\ref{thm:multiscale-parseval}): the increment identity and
  its telescoping sum were checked through $\ell=12$ for the Syracuse
  laws and for an independent compatible sequence of Dirichlet
  refinements. The largest absolute discrepancy was below
  $6\times10^{-14}$.
\item \emph{Logarithmic collision renormalization}
  (Theorem~\ref{thm:logarithmic-collision}): the collision mass from
  the original Syracuse recursion was compared with the independently
  evaluated spectral multiplier through $\ell=12$ at $q=3$, and
  through $\ell=5$ at $q=3,5,7,11,13,17,19$. These primes include
  several multiplicative orders and cases where $2$ is not a primitive
  root. The largest absolute discrepancy in the general test was below
  $2\times10^{-12}$ for collision masses as large as $10^4$; the
  logarithmic ordering was checked to be a permutation at every level.
\item \emph{Fresh-digit affine coupling}
  (Theorem~\ref{thm:fresh-digit-coupling}): exact enumeration for
  sibling gaps $2,4,8$, coarse precisions $1,2,3$, and fresh blocks
  through six digits reproduces $1-3^{-s}$ in total variation and
  $s\log3$ in mutual information. This also checks the orientation of
  the affine residue map and its carries.
\item \emph{$K_\ell$ recursion} (Empirical Result~\ref{thm:l2-finite}): the
  only initially available check was the hard-coded base case
  ($K_1=5/3$), which does not exercise the recursive machinery; we
  additionally validated the full distribution $\mu_\ell$, bin by bin,
  against direct Monte Carlo sampling of the primary definition
  \eqref{eq:syracuse-def} at $\ell=3,4$, within Monte Carlo noise.
\item \emph{$L^p$ collision spectrum}
  (Empirical Result~\ref{thm:lp-spectrum}): its $p=2$ column was
  compared level by level with the independently maintained
  $K_\ell$ recursion, agreeing to floating-point precision.
\item \emph{Microcanonical multiplicities}
  (Empirical Results~\ref{thm:microcanonical-finite} and
  \ref{thm:microcanonical-fourier}): residue sums were checked against
  independent bounded-composition coefficients at every level. The
  canonical mixture \eqref{eq:microcanonical} was compared bin by bin
  with the Syracuse recursion through $\ell=4$, with maximum
  discrepancy $2.1\times10^{-17}$. A count-valued implementation
  through $\ell=12$ and a Boolean implementation through $\ell=16$
  produced identical support thresholds on their common range.
\item \emph{Fixed-precision ensemble projection}
  (Empirical Result~\ref{thm:fixed-precision-finite}): each count row
  was projected directly modulo $3^r$ and compared in total variation
  with $\mu_r$ from the independently maintained Syracuse recursion.
  The comparison was run for every $r\le3$ and $r\le\ell\le12$ at the
  central and upper-window costs, then through full precision at
  $\ell=13$. At every level the code also checks the pointwise
  domination in Theorem~\ref{thm:ensemble-divergence}, using the exact
  bounded-composition count and folded-cost normalization.
\item \emph{Linear-block nonequivalence}
  (Theorem~\ref{thm:linear-block-nonequivalence}): exact folded-cost
  coefficients through $\ell=1000$ converge to the stated Gaussian
  total-variation limits. For $\rho=1/4,1/2,3/4$ and $u=0$, the limits
  are $0.0694909$, $0.1660641$, and $0.3226746$; the finite values at
  $\ell=1000$ are $0.0699624$, $0.1668085$, and $0.3241100$.
  Both finite distributions are normalized independently. The
  computation checks the constants; the proof uses the lattice local
  central limit theorem.
\item \emph{Jensen-identity Monte Carlo} (Theorem~\ref{thm:jensen}):
  the identity is exact and independent of simulation; the Monte Carlo
  estimate ($8\times10^6$ samples) serves only as a sanity check on
  the analytic value, and matches it to four significant figures.
\item \emph{Wirsching 2003 moment recursion} (Empirical Result~\ref{thm:conjecture3}):
  validated by an independent fixed-point check on an independent
  grid, agreeing to the grid's resolution ($\approx 6$ digits), and by
  confirming the functional equation~(7.12) of \cite{Wirsching2003}
  numerically up to $\ell=200$ using an analytic (not numerically
  differenced) logarithmic derivative.
\item \emph{Chang residue-balance test} (\S\ref{sec:context}):
  validated by confirming that no burst-ending event ever produced a
  residue outside the two predicted values ($9,25\bmod32$) in either
  the ensemble or single-orbit experiments, matching the source
  paper's structural proposition exactly.
\end{itemize}

\begin{thebibliography}{99}

\bibitem{Tao2022} T.\ Tao, \emph{Almost all orbits of the Collatz map
  attain almost bounded values}, Forum of Mathematics, Pi \textbf{10}
  (2022), e12. arXiv:1909.03562.

\bibitem{Tao2011} T.\ Tao, \emph{The Collatz conjecture, Littlewood--Offord
  theory, and powers of 2 and 3}, blog post,
  \url{https://terrytao.wordpress.com/2011/08/25/}, 2011.

\bibitem{Tao2020} T.\ Tao, \emph{Equidistribution of Syracuse random
  variables and density of Collatz preimages}, blog post,
  \url{https://terrytao.wordpress.com/2020/01/25/}, 2020.

\bibitem{Wirsching1998} G.\ J.\ Wirsching, \emph{The Dynamical System
  Generated by the $3n+1$ Function}, Lecture Notes in Mathematics
  \textbf{1681}, Springer, 1998.

\bibitem{Wirsching2003} G.\ J.\ Wirsching, \emph{On the problem of
  positive predecessor density in $3n+1$ dynamics}, Discrete and
  Continuous Dynamical Systems \textbf{9}(3) (2003), 771--787.
  DOI: 10.3934/dcds.2003.9.771.

\bibitem{LagariasBibliographyII} J.\ C.\ Lagarias, \emph{The $3x+1$
  problem: an annotated bibliography, II (2000--2009)},
  arXiv:math/0608208v6, 2012.

\bibitem{DragicevicEtAl2018} D.\ Dragi\v{c}evi\'c, G.\ Froyland,
  C.\ Gonz\'alez-Tokman, and S.\ Vaienti, \emph{A spectral approach
  for quenched limit theorems for random expanding dynamical systems},
  Communications in Mathematical Physics \textbf{360}(3) (2018),
  1121--1187. DOI: 10.1007/s00220-017-3083-7.

\bibitem{Hafouta2020} Y.\ Hafouta, \emph{Limit theorems for some
  time-dependent expanding dynamical systems}, Nonlinearity
  \textbf{33} (2020), 6421--6460.
  DOI: 10.1088/1361-6544/aba5e7.

\bibitem{KontorovichLagarias2009} A.\ V.\ Kontorovich and J.\ C.\
  Lagarias, \emph{Stochastic models for the $3x+1$ and $5x+1$
  problems}, in: \emph{The Ultimate Challenge: the $3x+1$ Problem}
  (J.\ C.\ Lagarias, ed.), American Mathematical Society, 2010.

\bibitem{AldousBandyopadhyay2005} D.\ Aldous and A.\ Bandyopadhyay,
  \emph{A survey of max-type recursive distributional equations},
  Annals of Applied Probability \textbf{15}(2) (2005), 1047--1110.

\bibitem{GoncalvesGreenfeldMadrid2022} F.\ Gon\c{c}alves, R.\
  Greenfeld, and J.\ Madrid, \emph{Generalized Collatz maps with almost
  bounded orbits}, arXiv:2111.06170, 2022.

\bibitem{DurrettLiggett1983} R.\ Durrett and T.\ M.\ Liggett,
  \emph{Fixed points of the smoothing transformation}, Zeitschrift für
  Wahrscheinlichkeitstheorie und verwandte Gebiete \textbf{64}(3)
  (1983), 275--301.

\bibitem{AlsmeyerBigginsMeiners2012} G.\ Alsmeyer, J.\ D.\ Biggins, and
  M.\ Meiners, \emph{The functional equation of the smoothing
  transform}, Annals of Probability \textbf{40}(5) (2012), 2069--2105.

\bibitem{JelenkovicOlveraCravioto2012} P.\ R.\ Jelenkovi\'c and
  M.\ Olvera-Cravioto, \emph{Implicit renewal theorem for trees with
  general weights}, Stochastic Processes and their Applications
  \textbf{122}(9) (2012), 3209--3238.

\bibitem{VillemonaisZalduendo2025} D.\ Villemonais and N.\ Zalduendo,
  \emph{Convergence of weighted branching processes}, arXiv:2512.07653
  (2025).

\bibitem{KoleskoMentemeier2015} K.\ Kolesko and S.\ Mentemeier,
  \emph{Fixed points of the multivariate smoothing transform: the
  critical case}, arXiv:1409.7220.

\bibitem{Goldie1991} C.\ M.\ Goldie, \emph{Implicit renewal theory and
  tails of solutions of random equations}, Annals of Applied
  Probability \textbf{1}(1) (1991), 126--166.

\bibitem{BuraczewskiDamekMikosch2016} D.\ Buraczewski, E.\ Damek, and
  T.\ Mikosch, \emph{Stochastic Models with Power-Law Tails: The
  Equation $X=AX+B$}, Springer, 2016.

\bibitem{KrasikovLagarias2003} I.\ Krasikov and J.\ C.\ Lagarias,
  \emph{Bounds for the $3x+1$ problem using difference inequalities},
  Acta Arithmetica \textbf{109} (2003), 237--258.

\bibitem{Hooley1967} C.\ Hooley, \emph{On Artin's conjecture}, Journal
  f\"ur die reine und angewandte Mathematik \textbf{225} (1967),
  209--220.

\bibitem{Biggins1992} J.\ D.\ Biggins, \emph{Uniform convergence of
  martingales in the branching random walk}, Annals of Probability
  \textbf{20}(1) (1992), 137--151.

\bibitem{Nerman1981} O.\ Nerman, \emph{On the convergence of
  supercritical general (C-M-J) branching processes}, Zeitschrift
  f\"ur Wahrscheinlichkeitstheorie und verwandte Gebiete
  \textbf{57}(3) (1981), 365--395.

\bibitem{ApplegateLagarias1995I} D.\ Applegate and J.\ C.\ Lagarias,
  \emph{Density bounds for the $3x+1$ problem I: tree-search method},
  Mathematics of Computation \textbf{64} (1995), 411--426.

\bibitem{ApplegateLagarias1995II} D.\ Applegate and J.\ C.\ Lagarias,
  \emph{Density bounds for the $3x+1$ problem II: Krasikov
  inequalities}, Mathematics of Computation \textbf{64} (1995),
  427--438.

\bibitem{Liu2000} Q.\ Liu, \emph{On generalized multiplicative
  cascades}, Stochastic Processes and their Applications
  \textbf{86}(2) (2000), 263--286.

\bibitem{BakerBanaji2026} S.\ Baker and A.\ Banaji, \emph{Self-similar
  and self-conformal measures with slow Fourier decay},
  arXiv:2602.05593, 2026.

\bibitem{Bremont} J.\ Br\'emont, \emph{Self-similar measures and the
  Rajchman property}, Annales Henri Lebesgue \textbf{4} (2021),
  973--1004. arXiv:1910.03463.

\bibitem{LiSahlsten2022} J.\ Li and T.\ Sahlsten, \emph{Trigonometric
  series and self-similar sets}, Journal of the European Mathematical
  Society \textbf{24}(1) (2022), 341--368.

\bibitem{Solomyak1995} B.\ Solomyak, \emph{On the random series
  $\sum\pm\lambda^n$ (an Erd\H{o}s problem)}, Annals of Mathematics
  \textbf{142}(3) (1995), 611--625.

\bibitem{VarjuYu2022} P.\ Varj\'u and H.\ Yu, \emph{Fourier decay of
  self-similar measures and self-similar sets of uniqueness}, Analysis
  \& PDE \textbf{15}(3) (2022), 843--858. arXiv:2004.09358.

\bibitem{Spiegelhofer2021} L.\ Spiegelhofer, \emph{Collisions of digit
  sums in bases 2 and 3}, Israel Journal of Mathematics \textbf{258}
  (2023), 475--502. arXiv:2105.11173.

\bibitem{Chang2026} E.\ Y.\ Chang, \emph{A structural reduction of the
  Collatz conjecture to one-bit orbit mixing}, arXiv:2603.25753, 2026.

\bibitem{ChangCompanion2026} E.\ Y.\ Chang, \emph{Exploring Collatz
  dynamics with human--LLM collaboration}, arXiv:2603.11066, 2026.

\bibitem{Halasz1977} G.\ Hal\'asz, \emph{Estimates for the
  concentration function of combinatorial number theory and
  probability}, Periodica Mathematica Hungarica \textbf{8}(3--4)
  (1977), 197--211.

\bibitem{TaoVu2009} T.\ Tao and V.\ Vu, \emph{Inverse Littlewood--Offord
  theorems and the condition number of random discrete matrices},
  Annals of Mathematics \textbf{169} (2009), 595--632.

\bibitem{NguyenVu2011} H.\ Nguyen and V.\ Vu, \emph{Optimal inverse
  Littlewood--Offord theorems}, Advances in Mathematics \textbf{226}
  (2011), 5298--5319.

\bibitem{Mahler1968} K.\ Mahler, \emph{An unsolved problem on the powers
  of $3/2$}, Journal of the Australian Mathematical Society
  \textbf{8}(2) (1968), 313--321.

\bibitem{FlattoLagariasPollington1995} L.\ Flatto, J.\ C.\ Lagarias,
  and A.\ D.\ Pollington, \emph{On the range of fractional parts
  $\{\xi(p/q)^n\}$}, Acta Arithmetica \textbf{70}(2) (1995), 125--147.

\bibitem{Erdos1979} P.\ Erd\H{o}s, \emph{Some unconventional problems
  in number theory}, Mathematics Magazine \textbf{52}(2) (1979),
  67--70.

\bibitem{Lagarias2009} J.\ C.\ Lagarias, \emph{Ternary expansions of
  powers of 2}, Journal of the London Mathematical Society
  \textbf{79}(3) (2009), 562--588.

\end{thebibliography}

\end{document}
